% 開発記録（ゼミ報告・将来の研究論文用）
% 物理モデル自動構築（GNN）のための数式抽出パイプライン
\documentclass[11pt,a4paper]{jsarticle}
\usepackage[T1]{fontenc}
\usepackage[utf8]{inputenc}
\usepackage{amsmath,amssymb}
\usepackage{url}
\usepackage[dvipdfmx]{hyperref}
\usepackage{enumitem}
\usepackage{geometry}
\geometry{margin=2.5cm}

\title{開発記録：科学文献からの数式抽出パイプライン}
\author{AutoPMob プロジェクト}
\date{更新日：\today}

\begin{document}
\maketitle
\tableofcontents
\clearpage

%----------------------------------------------------------------------
\section{概要}
%----------------------------------------------------------------------
本ドキュメントは、GNN（グラフニューラルネットワーク）を用いた物理モデル自動構築の研究の一環として開発している、科学文献PDFから数式と変数定義を抽出するパイプラインの開発記録である。
記録内容は研究室ゼミの報告および将来の研究論文の素材として利用する。

\textbf{パイプラインの目的}
\begin{itemize}[nosep]
  \item 科学文献のPDFを入力とし、番号付き数式および付録の数式を LaTeX 形式で抽出する。
  \item 各数式について変数定義（variable definitions）と物理的意味（context）を周辺テキストから抽出する。
  \item 抽出結果を JSON 形式で保存し、物理モデル自動生成モデルの訓練データとして利用する。
\end{itemize}

\textbf{技術スタック}
\begin{itemize}[nosep]
  \item API: Google Gemini API（\texttt{gemini-2.5-flash}）
  \item SDK: \texttt{google-genai}（PDF 直接アップロード、Structured Output 対応）
  \item スキーマ検証: Pydantic
  \item 出力: \texttt{extracted\_equations.json}（および表記正規化版 \texttt{extracted\_equations\_normalized.json}）
\end{itemize}

%----------------------------------------------------------------------
\section{遭遇した問題と解決策}
%----------------------------------------------------------------------

\subsection{API モデル 404（モデル未提供）}
\textbf{現象}
\begin{itemize}[nosep]
  \item \texttt{gemini-1.5-pro} を指定して \texttt{generateContent} を呼ぶと、API から \texttt{404 NOT\_FOUND} が返る。
  \item エラーメッセージ: ``models/gemini-1.5-pro is not found for API version v1beta''.
\end{itemize}

\textbf{原因}
\begin{itemize}[nosep]
  \item 利用環境（API version v1beta）では \texttt{gemini-1.5-pro} が提供されていない、またはモデルIDが変更されている。
\end{itemize}

\textbf{対応}
\begin{itemize}[nosep]
  \item 呼び出しモデルを \texttt{gemini-2.0-flash} に変更。のちに無料枠の都合で \texttt{gemini-2.5-flash} に変更。
  \item 利用可能なモデル一覧は \url{https://ai.google.dev/gemini-api/docs/models} 等で確認する。
\end{itemize}

%----------------------------------------------------------------------
\subsection{レートリミット（429 RESOURCE\_EXHAUSTED）}
\textbf{現象}
\begin{itemize}[nosep]
  \item リクエスト直後に \texttt{429 RESOURCE\_EXHAUSTED} が返る。
  \item メッセージに ``Quota exceeded for metric: ... generate\_content\_free\_tier\_..., limit: 0'' とあり、無料枠のリクエスト数・トークン数が 0 に制限されている。
\end{itemize}

\textbf{原因}
\begin{itemize}[nosep]
  \item 無料ティアのクォータが枯渇している、または当該プロジェクトでは \texttt{gemini-2.0-flash} の無料枠が 0 に設定されている。
\end{itemize}

\textbf{対応}
\begin{itemize}[nosep]
  \item 利用可能なモデルに変更：\texttt{gemini-2.5-flash} では「1分あたり5回（5 RPM）」の枠が利用可能であったため、モデルを \texttt{gemini-2.5-flash} に統一。
  \item 5 RPM を超えないよう、\texttt{generate\_content} の呼び出し前に「前回呼び出しから最低12秒（60/5）経過するまで待機」する処理を追加。
  \item リトライ回数・待機時間は当初「12回・最大300秒」としていたが、ユーザ体験を考慮して「5回・初回15秒・最大60秒」に短縮（レートリミット時のみリトライ）。
\end{itemize}

%----------------------------------------------------------------------
\subsection{数式内の変数表記の揺れ（$c_w$ と $c_p$）}
\textbf{現象}
\begin{itemize}[nosep]
  \item 論文では5番・6番の式で「水の比熱」を明示するため $c_w$ を用いているが、抽出結果ではすべて $c_p$ になっていた（eq\_5, eq\_6）。
\end{itemize}

\textbf{原因（推定）}
\begin{itemize}[nosep]
  \item \textbf{LLM の記号正規化}：熱力学では比熱の記号として $c_p$ が一般的なため、モデルが $c_w$ を $c_p$ に「統一」して出力した可能性。
  \item \textbf{PDF の見た目}：フォントによっては下付きの $w$ と $p$ が似て見え、誤読した可能性。
  \item \textbf{意味の同一視}：「水で満たされている」という文脈で $c_w$ と $c_p$ を同一視し、$c_p$ に揃えた可能性。
\end{itemize}

\textbf{対応}
\begin{itemize}[nosep]
  \item プロンプトを強化：「変数記号・添字は文献の表記をそのまま抽出する。$c_w$ を $c_p$ に変更しない」旨を明記。
  \item 訓練用には「同じ物理意味の変数は正規化する」方針とし、\texttt{config/notation\_map.json} と \texttt{normalize\_notation.py} を用意。$c_w$/Cp/cp 等を $c_p$ に統一する後処理を導入。抽出はできるだけ原文のまま、訓練前に正規化する二段構えにした。
\end{itemize}

%----------------------------------------------------------------------
\subsection{変数表記の統一（訓練用正規化）}
\textbf{要望}
\begin{itemize}[nosep]
  \item 物理モデル自動生成モデルの訓練のため、同じ物理意味を持つ変数の表記をすべて統一したい（モデル誤読・論文側の表記揺れのいずれも吸収）。
\end{itemize}

\textbf{対応}
\begin{itemize}[nosep]
  \item \texttt{config/notation\_map.json} で「正規表記（canonical）」と「揺れ（alternatives）」を定義。
  \item \texttt{normalize\_notation.py} で、抽出済み JSON の \texttt{equation} および \texttt{variables} のキー・説明文に置換を適用し、\texttt{*\_normalized.json} を出力。
  \item 抽出パイプラインに \texttt{--normalize} オプションを追加し、抽出直後に正規化版を保存できるようにした。
\end{itemize}

%----------------------------------------------------------------------
\section{技術メモ}
%----------------------------------------------------------------------
\begin{itemize}[nosep]
  \item \textbf{Structured Output}：Pydantic の \texttt{model\_json\_schema()} で配列スキーマを組み、\texttt{response\_mime\_type: application/json} と \texttt{response\_json\_schema} を渡して JSON 配列を取得。
  \item \textbf{ファイルアップロード}：\texttt{client.files.upload()} で PDF をアップロードし、\texttt{state} が \texttt{ACTIVE} になるまでポーリング。\texttt{Part.from\_uri()} でコンテンツに含める。
  \item \textbf{環境変数}：\texttt{GEMINI\_API\_KEY} または \texttt{GOOGLE\_API\_KEY}。オプションで \texttt{EXTRACT\_EQUATIONS\_MIN\_INTERVAL}（5 RPM 用）、\texttt{EXTRACT\_EQUATIONS\_MAX\_RETRIES} 等。
  \item \textbf{開発記録の自動追記と GitHub 同期}：Cursor のルール（\texttt{.cursor/rules/development-log.mdc}）で、\textbf{全ての変動}（機能追加・新スクリプト・バグ修正・設定変更等）の内容と理由を \texttt{development\_log.tex} の「更新履歴」に追記するよう指示。あわせて、変更後は必ず \texttt{git add}、\texttt{git commit}、\texttt{git push} を実行し、GitHub へ同期する運用とする。
  \item \textbf{コミット時の自動 push}：\texttt{.git/hooks/post-commit} で \texttt{git commit} のたびに \texttt{git push} を実行するため、コミットすればリモートへ反映される。
  \item \textbf{統合パイプライン \texttt{build\_equation\_graph.py}}：\texttt{input\_pdfs/} 内の「未抽出」PDF のみ一括処理（処理済みは \texttt{processed\_pdfs.json} で管理、1ファイルごとに 15\,s 待機で 5\,RPM 対策）。LLM 知識から化学工学基礎式を 20 個生成（Handbook）。抽出・生成結果を単一の \texttt{unified\_equations.json} に統合（既存時は追記、\texttt{source\_id}+\texttt{eq\_id} で重複排除）。\texttt{--mode pdf|handbook|all}、\texttt{--processed-list} で処理済み一覧を指定可能。
  \item \textbf{API モデル}：当初 \texttt{gemini-1.5-pro}、のち \texttt{gemini-2.5-flash}（5 RPM）、RPD 制限対策で \texttt{gemini-2.0-flash} に変更後、高性能化のため \texttt{gemini-2.5-pro} に統一。
  \item \textbf{デフォルト API キー}：環境変数未設定時に限り、開発用デフォルトキーを使用するように変更（\texttt{extract\_equations.py}、\texttt{build\_equation\_graph.py}、\texttt{test\_api.py}）。
  \item \textbf{論文メタデータ取得 \texttt{fetch\_papers.py}}：Semantic Scholar API でクエリ A/B 各 50 件のメタデータを取得し、\texttt{paper\_dataset\_links.json} に保存。オープンアクセス PDF URL を \texttt{has\_open\_access\_pdf}、\texttt{open\_access\_pdf\_url} で明示。429 時はリトライと \texttt{SEMANTIC\_SCHOLAR\_API\_KEY} 対応。
  \item \textbf{グラフ前処理 \texttt{build\_graph\_data.py}}：\texttt{unified\_equations.json} から\textbf{式--変数二部グラフ}の \texttt{Data} を構築し \texttt{equation\_graph.pt} に保存（詳細は\nameref{sec:file-roles}）。
  \item \textbf{訓練ケース生成 \texttt{generate\_training\_cases.py}}：LLM で約150件のコアケースを生成（\texttt{context}、\texttt{input\_variables}、\texttt{output\_variables}、\texttt{correct\_model\_ids}）。正解モデルは「入力が与えられれば出力を解ける」「十分かつ最小限の数式群」となるようプロンプトで制約。Python で入出力スワップ・ランダム入出力・context 言い換えにより 1 コアから 5--6 バリアントを生成し、900件以上を \texttt{training\_cases.json} に保存。
  \item \textbf{可視化・分析 \texttt{analyze\_dataset.py}}：数式グラフ・訓練ケースの統計と可視化。ドメイン分布、次数分布、埋め込みの PCA（ドメイン色分け）、変数共起ネットワーク、訓練ケースの variant 分布・正解モデル数分布。論文付録向け \texttt{dataset\_report.txt} と \texttt{figures/*.png} を出力。
\end{itemize}

%----------------------------------------------------------------------
\section{各ファイル・ディレクトリの役割}
\label{sec:file-roles}
%----------------------------------------------------------------------
リポジトリ内の主要なパスと、その責務の一覧である（\texttt{.gitignore} 対象や個人環境依存は除く）。

\subsection{プロジェクトルート（実行スクリプト・メタファイル）}
\begin{itemize}[nosep]
  \item \texttt{README.md}：環境構築、代表的なコマンド、リポジトリ構成の入口。
  \item \texttt{CLAUDE.md}：別ツール・別環境で開発を続けるときの引き継ぎ要約（目的・データ・実験結論・スクリプト対応）。
  \item \texttt{research\_context.md}：研究メモ用（\texttt{CLAUDE.md} へのポインタを含む）。
  \item \texttt{requirements.txt}：Python 依存パッケージの一覧。
  \item \texttt{extract\_equations.py}：単一 PDF を Gemini API に渡し、数式・変数定義・周辺文脈を JSON で抽出する。
  \item \texttt{normalize\_notation.py}：抽出済み JSON の式文字列・変数キー等を \texttt{config/notation\_map.json} に従い訓練向けに正規化する。
  \item \texttt{build\_equation\_graph.py}：\texttt{input\_pdfs/} の未処理 PDF を一括抽出し、必要に応じ Handbook 式を LLM 生成し、\texttt{unified\_equations.json} に統合する（スキップ管理は \texttt{processed\_pdfs.json}）。
  \item \texttt{fetch\_papers.py}：Semantic Scholar API で論文メタデータを取得し \texttt{paper\_dataset\_links.json} に保存する。
  \item \texttt{build\_graph\_data.py}：\texttt{unified\_equations.json} から式ノード・変数ノードの二部グラフを構築し、各ノードに CodeBERT [CLS]（768次元、変数側は含有式の平均）を付与して \texttt{equation\_graph.pt}（PyTorch Geometric）を出力する。
  \item \texttt{generate\_training\_cases.py}：LLM でコアケースを生成し、入出力スワップ等でバリアントを増やして \texttt{training\_cases.json} を構築する。
  \item \texttt{validate\_dataset.py}：式のグローバル ID 参照、ケースと式の変数整合、重複などを検査し \texttt{validate\_dataset\_report.json} を出力する。
  \item \texttt{add\_context\_paraphrases.py}：訓練ケースの文脈のみ言い換えたバリアントを \texttt{training\_cases.json} に追加する。
  \item \texttt{analyze\_dataset.py}：グラフ・訓練ケースの統計と図表を生成し、\texttt{dataset\_report.txt} と \texttt{figures/} に書き出す。
  \item \texttt{baseline\_retrieval\_eval.py}：TF-IDF、変数 Jaccard、混合ベースライン等の retrieval 評価ロジック（単体実行および実験ランナーから利用）。
  \item \texttt{run\_retrieval\_experiments.py}：複数の scoring 方式を同一の source\_id 分割で学習・評価し、結果を \texttt{experiments/} に JSON と Markdown で保存する。
  \item \texttt{train\_gnn\_retriever.py}：\texttt{equation\_graph.pt} 上で式側を GCN または GAT で更新し、クエリを SVD$\to$MLP で揃えて InfoNCE 型で学習する簡易スクリプト（レポートは \texttt{gnn\_train\_eval\_report.json}）。
  \item \texttt{train\_gnn\_refine\_svd.py}：式・クエリを同一 TF-IDF+SVD 空間に置き、式埋め込みを残差 GCN で更新する実験用スクリプト（\texttt{gnn\_refine\_svd\_report.json}）。
  \item \texttt{test\_api.py}：Gemini API の疎通確認用。
\end{itemize}

\subsection{設定・運用補助}
\begin{itemize}[nosep]
  \item \texttt{config/notation\_map.json}：同一物理量の記号揺れを正規表記に寄せるための対応表。
  \item \texttt{scripts/install-git-hooks.sh}：コミット後に \texttt{git push} する \texttt{post-commit} フックを入れるためのインストール用シェル。
  \item \texttt{.cursor/rules/development-log.mdc}：コード変更時に本 \texttt{development\_log.tex} を更新し GitHub に同期するよう指示するエディタ用ルール。
\end{itemize}

\subsection{データ・実験成果物（リポジトリに含まれる例）}
\begin{itemize}[nosep]
  \item \texttt{unified\_equations.json}：全ソースを統合した数式ライブラリ（グラフ構築・訓練ケースの参照元）。
  \item \texttt{training\_cases.json}：retrieval タスクの入力（文脈・I/O 変数）と正解式 ID 集合。
  \item \texttt{equation\_graph.pt}：PyG \texttt{Data}（\texttt{x}, 二部 \texttt{edge\_index}, \texttt{num\_equations}, \texttt{edge\_index\_eq\_eq} 等）。
  \item \texttt{processed\_pdfs.json}：\texttt{build\_equation\_graph.py} が既に処理した PDF ファイル名のリスト。
  \item \texttt{paper\_dataset\_links.json}：\texttt{fetch\_papers.py} の取得結果（PDF URL 等）。
  \item \texttt{extracted\_equations.json}：単一 PDF 抽出の典型出力（中間・サンプル用）。
  \item \texttt{experiments/retrieval\_experiments.json}：方式比較実験の分割・ハイパラ・指標の機械可読ログ。
  \item \texttt{experiments/retrieval\_experiments.md}：上記実験の表形式サマリ。
  \item \texttt{validate\_dataset\_report.json}：データセット検証の詳細結果。
  \item \texttt{dataset\_report.txt}：\texttt{analyze\_dataset.py} のテキスト要約。
  \item \texttt{gnn\_train\_eval\_report.json} / \texttt{gnn\_refine\_svd\_report.json} / \texttt{baseline\_eval\_report.json}：各スクリプトの評価サマリ（実行時に更新）。
  \item \texttt{input\_pdfs/}：一括抽出の入力 PDF を置くディレクトリ。
  \item \texttt{figures/}：分析スクリプトが出力する図（PNG 等）。
\end{itemize}

\subsection{ドキュメント}
\begin{itemize}[nosep]
  \item \texttt{docs/development\_log.tex}：本ドキュメント（開発記録の原本）。
  \item \texttt{docs/README.md}：\texttt{development\_log.tex} を PDF にビルドする手順。
  \item \texttt{docs/GITHUB\_SYNC.md}：リモート同期に関する補足メモ。
\end{itemize}

%----------------------------------------------------------------------
\section{データセット定義と評価（次の開発者向けガイド）}
%----------------------------------------------------------------------

\subsection{データ構造（Equation と Training Case）}
\textbf{Equation（\texttt{unified\_equations.json}）}
\begin{itemize}[nosep]
  \item 1要素が1つの数式（ノード）で、主なキーは \texttt{source\_id}, \texttt{eq\_id}, \texttt{equation}, \texttt{variables}, \texttt{context\_text}, \texttt{domain}。
  \item \texttt{variables} は辞書で、キーが変数シンボル（例: \texttt{C\_A}）、値が物理意味の説明文（例: Concentration of component A）。
  \item \textbf{グローバルIDの注意}: \texttt{eq\_id} は文書内ローカルのことがあるため、訓練や評価では \textbf{必ず} \texttt{<source\_id>\_\_<eq\_id>} 形式（例: \texttt{handbook\_chemeng\_\_eq\_1}）を「式ID」として扱う。
\end{itemize}

\textbf{Training Case（\texttt{training\_cases.json}）}
\begin{itemize}[nosep]
  \item 1要素が1つの推論タスクで、主なキーは \texttt{case\_id}, \texttt{variant\_type}, \texttt{context}, \texttt{input\_variables}, \texttt{output\_variables}, \texttt{correct\_model\_ids}。
  \item \texttt{correct\_model\_ids} は「このケースの正解となる数式集合（1本以上）」であり、各要素は \textbf{必ず} \texttt{<source\_id>\_\_<eq\_id>} 形式。
  \item \textbf{variant\_type}: \texttt{original}, \texttt{swap\_io}, \texttt{random\_io\_from\_models}, \texttt{context\_paraphrased}（I/Oと正解モデルを保持し文脈だけ言い換え）。
\end{itemize}

\subsection{タスク定義（Equation Retrieval / Model Selection）}
本研究の基本タスクは、与えられたケース $(\texttt{context}, \texttt{input\_variables}, \texttt{output\_variables})$ から、式集合ライブラリ（全式）内の \texttt{correct\_model\_ids} を上位にランキングする \textbf{retrieval} として定義する。
各ケースは正解が1本のことも複数本のこともある（multi-positive）。

\subsection{なぜ MRR を使うか（採用理由）}
本タスクでは「正解が上位に出てくること」が重要であり、特に \textbf{1本目の正解}が何位に現れるかがユーザ体験（候補式提示）に直結する。
また、ケースによって正解式が複数あるため、分類精度（accuracy）のような単一ラベル前提の指標は不適切である。
そのため、ランキングの早期精度を評価できる \textbf{MRR（Mean Reciprocal Rank）} と \textbf{Recall@K} を主要指標として採用した。

\subsection{MRR と Recall@K の計算方法}
各ケース $i$ について、モデルが全式をスコアリングし降順に並べた順位列を考える。
正解集合を $G_i$（\texttt{correct\_model\_ids} に対応するノード集合）とする。
\begin{itemize}[nosep]
  \item \textbf{BestRank}：$r_i=\min\{ \text{rank}(e)\mid e\in G_i\}$（正解集合のうち最も上位に来た式の順位）
  \item \textbf{Reciprocal Rank}：$RR_i=\frac{1}{r_i}$
  \item \textbf{MRR}：$\mathrm{MRR}=\frac{1}{N}\sum_{i=1}^{N} RR_i$
  \item \textbf{Recall@K}：$\mathrm{Recall@K}=\frac{1}{N}\sum_{i=1}^{N}\mathbb{I}[r_i\le K]$
\end{itemize}
ここで $N$ は評価対象ケース数。正解が複数ある場合でも、最初に当たる正解の順位（BestRank）で評価する。

\subsection{データ分割（リーク防止のための paper/source\_id split）}
同一文書から抽出された式や文脈が train と test に混ざると、文字列一致により過大評価（リーク）になる。
そのため、分割は \textbf{source\_id（論文/教科書）単位}で行い、case の \texttt{correct\_model\_ids} から source\_id を推定して train/val/test に割り当てる。
実装例は \texttt{run\_retrieval\_experiments.py}（\texttt{experiments/retrieval\_experiments.json} に split を保存）。

\subsection{再現可能な実験手順（何をどう動かすか）}
\begin{itemize}[nosep]
  \item \textbf{健全性チェック}：\texttt{python validate\_dataset.py}（式ID参照、I/O変数整合、重複など。レポートは \texttt{validate\_dataset\_report.json}）。
  \item \textbf{訓練ケース増強（contextのみ言い換え）}：\texttt{python add\_context\_paraphrases.py}（\texttt{training\_cases.json} を更新）。
  \item \textbf{統計と図の更新}：\texttt{python analyze\_dataset.py}（\texttt{dataset\_report.txt} と \texttt{figures/}）。
  \item \textbf{方式比較とハイパラ探索}：\texttt{python run\_retrieval\_experiments.py}（結果は \texttt{experiments/retrieval\_experiments.md}）。
\end{itemize}

\subsection{現時点の結論（次に何をすべきか）}
現状では、学習を伴うGNN/MLP系よりも、\textbf{TF-IDF（文脈）+ 変数集合Jaccard（I/O一致）}を混合した \texttt{baseline\_mix} が最良である。
次の開発では、(1) negative sampling を hard negative 化、(2) 変数をノードとして明示した二部グラフ化、(3) 「正解式集合の整合性（最小十分性）」を利用した学習目的の改善、などによりベースライン超えを狙う。

%----------------------------------------------------------------------
\section{GNNの仕組みと埋め込み（次の開発者向け）}
%----------------------------------------------------------------------

\subsection{なぜ現状のグラフ構築ではGNNの強みを活かせないか（分析）}
GNNの強みは、\textbf{タスクに効く構造をグラフに載せ、その上で局所的な集約をしてノード表現を磨く}ことにある。当初の「式をノード、共通変数があればエッジ」という設計には、次のような限界がある。

\textbf{（1）エッジが極端に密になり、構造の弁別力が弱い}
\begin{itemize}[nosep]
  \item ノード数 $N$ に対し、変数を1つでも共有する式ペアに無向エッジを張ると、$t$, $T$, $V$, $C_A$ など頻出変数を含む式が多数つながる。実際のデータでは 975 ノードで約 16 万本のエッジとなり、1式あたり平均で非常に多くの式と隣接する。
  \item その結果、「同じ物理モデルに属する式」と「たまたま $t$ や $T$ を共有するだけの式」が区別されず、メッセージパッシングで無関係な情報が混ざり、ノード表現がシャープに磨かれない。
\end{itemize}

\textbf{（2）変数共有と retrieval の正解が対応していない}
\begin{itemize}[nosep]
  \item タスクは「与えられたケース（context + I/O変数）に合う式をランキングする」retrieval である。正解は「そのケースの入出力を満たす式集合」であり、変数名の共有は「候補になりうる」必要条件にすぎない。
  \item 多くの不正解の式も $t$ や $T$ を共有するため、グラフのつながりだけでは「このクエリに正解」というシグナルが弱く、GNN の集約が retrieval に直結しにくい。
\end{itemize}

\textbf{（3）変数がグラフ上に明示されていない}
\begin{itemize}[nosep]
  \item 変数はエッジの「条件」としてしか現れず、ノードとしては存在しない。そのため「この式はクエリの出力変数 $C_A$ を含む」といった情報を、メッセージパッシングで明示的に扱えない。のちにクエリ条件（I/O変数）に依存した重み付けやサブグラフを入れることも難しい。
\end{itemize}

\textbf{（4）ホモフィリー仮定が成り立ちにくい}
\begin{itemize}[nosep]
  \item 多くの GCN 系モデルは「つながっているノードは似ている」という仮定で学習する。ここでは「つながり＝1つでも変数名が被った」なので、ドメインが違っても $t$ だけ共有でつながる、といったノイズが多く、ホモフィリーが弱い。そのため、隣接ノードからの集約が「関連の強い式」に絞られず、GNN の利点が生かしづらい。
\end{itemize}

以上から、\textbf{「式＝ノード・共通変数＝エッジ」だけの同質グラフは、retrieval というタスクに対して GNN の強みを十分に引き出せる設計にはなっていない}、という結論になる。改善の方向性として、(i) 変数を明示した\textbf{式--変数二部グラフ}で「どの変数を介してつながるか」を構造に載せる、(ii) クエリの I/O 変数に応じたエッジ重みやサブグラフにする、(iii) 訓練ラベルを利用して「同じ正解に共起する式」同士をエッジで結ぶ、などが考えられる。

\subsection{グラフの構築（\texttt{build\_graph\_data.py} → \texttt{equation\_graph.pt}）【二部グラフ版】}
タスクに効く構造と局所集約を両立させるため、\textbf{式ノードと変数ノードの二部グラフ}を採用する。

\begin{itemize}[nosep]
  \item \textbf{ノード}：\textbf{式ノード}（0 から $N_{\mathrm{eq}}-1$）と \textbf{変数ノード}（$N_{\mathrm{eq}}$ から $N_{\mathrm{eq}}+N_{\mathrm{var}}-1$）。式ノードは \texttt{unified\_equations.json} の並びと一致。変数ノードは全式から出現する変数名のユニーク集合に 1 対 1 で対応する。
  \item \textbf{エッジ}：式 $i$ が変数 $v$ を含むときのみ、無向エッジ $(i,\; N_{\mathrm{eq}}+j)$ と $(N_{\mathrm{eq}}+j,\; i)$ を張る（$j$ は変数 $v$ の ID）。式同士は直接つながらず、\textbf{変数ノードを介した 2 ホップ}で「同じ変数を含む式」に情報が伝わる。これにより集約が「どの変数を共有しているか」に沿った形になり、局所性も保たれやすい。
  \item \textbf{ノード特徴 \texttt{x}}：式ノードには従来どおり \textbf{CodeBERT} の [CLS] ベクトル（768 次元）。変数ノードには、その変数を含む式の CodeBERT ベクトルの平均を 768 次元特徴として与える（学習可能な変数埋め込みにすることも可能）。\texttt{x} の形状は $(N_{\mathrm{eq}}+N_{\mathrm{var}},\; 768)$。
  \item \textbf{保存内容}：\texttt{Data(x=x, edge\_index=edge\_index, num\_equations=}$N_{\mathrm{eq}}$\texttt{)} に加え、後方互換のため式どうしの隣接 \texttt{edge\_index\_eq\_eq}（0..$N_{\mathrm{eq}}-1$ のみのインデックス）も保存する。学習・評価では式埋め込みは GCN 出力の \textbf{先頭 $N_{\mathrm{eq}}$ 行}のみを用いる。
\end{itemize}

\subsection{式の埋め込みの2系統}
Retrieval 実験では、式の表現に次の2系統がある（スクリプトによってどちらを使うかが異なる）。

\textbf{系統A：CodeBERT ノード特徴（\texttt{equation\_graph.pt} の \texttt{x}）}
\begin{itemize}[nosep]
  \item 上記のとおり、事前学習済み CodeBERT で式+文脈をエンコードした 768 次元ベクトル。学習時には L2 正規化してコサイン類似度＝内積として使う。
  \item \texttt{train\_gnn\_retriever.py} の「式側は \texttt{x} をそのまま（または GCN で更新）」や、\texttt{run\_retrieval\_experiments.py} の \texttt{q\_mlp\_to\_x} で使用。クエリ側は SVD 空間のベクトルを MLP で 768 次元に写す方式。
\end{itemize}

\textbf{系統B：TF-IDF + TruncatedSVD（同一空間でクエリと式を比較）}
\begin{itemize}[nosep]
  \item 式テキスト：\texttt{context\_text}, \texttt{equation}, \texttt{domain} を連結した文字列を TF-IDF（\texttt{TfidfVectorizer}、lowercase, max\_features=50000, ngram\_range=(1,2)）でベクトル化し、TruncatedSVD で 256 次元に圧縮。L2 正規化したものを式埋め込み $E_0$ とする。
  \item クエリ（ケース）テキスト：\texttt{context} と \texttt{INPUT <input\_variables>}, \texttt{OUTPUT <output\_variables>} を連結し、同じ TF-IDF と SVD で 256 次元に写し L2 正規化したものをクエリ埋め込み $Q$ とする。
  \item スコアは $E_0^\top Q$（内積＝コサイン類似度）。\texttt{svd\_cos} ベースラインや \texttt{train\_gnn\_refine\_svd.py} の入力はこの空間。
\end{itemize}

\subsection{GNN の役割（2種類の使い方）}
\textbf{パターン1：式特徴を CodeBERT \texttt{x} のまま、GCN で更新（\texttt{train\_gnn\_retriever.py}）}
\begin{itemize}[nosep]
  \item 式埋め込みは $z = \mathrm{GCN}(x, \mathrm{edge\_index})$。2層 GCN（\texttt{GCNConv(in\_dim, 768)}, ReLU, Dropout, \texttt{GCNConv(768, 768)}）の出力を L2 正規化。クエリは SVD ベクトルを \texttt{QueryEncoder}（MLP）で 768 次元に写し、$z^\top q$ でスコア。InfoNCE 風の負例サンプリングで学習。
  \item 現状の実験では、クエリと式が別空間（SVD vs CodeBERT）から来るため学習が難しく、このパターン単体ではベースラインを下回っている。
\end{itemize}

\textbf{パターン2：SVD 空間で式埋め込みを GCN で残差更新（\texttt{train\_gnn\_refine\_svd.py}, \texttt{gnn\_refine\_svd\_residual}）}
\begin{itemize}[nosep]
  \item 式の初期埋め込みは系統Bの $E_0$（256次元、L2正規化済み）。これを 2 層 GCN の入力とし、$h = \mathrm{GCN}(E_0)$ を計算したあと、\textbf{残差混合} $E = E_0 + \alpha \cdot h$ で更新（$\alpha$ は学習可能スカラー、初期値 0 でクリップ）。最終的に $E$ を L2 正規化して式埋め込みとする。クエリは同じ SVD 空間の $Q$ のまま、スコアは $E^\top Q$。
  \item $\alpha=0$ のときは GCN を通さないので SVD コサインと同一。グラフ構造（変数共有）の情報を少しずつ取り込む設計。学習はサンプルされた負例に対する InfoNCE 的損失。
\end{itemize}

\subsection{スコアリングと学習}
\begin{itemize}[nosep]
  \item いずれの方式でも、式埋め込みとクエリ埋め込みは L2 正規化してから内積をとるため、スコアは \textbf{コサイン類似度} に等しい。
  \item 学習時：各訓練ケースについて正解式を1本選び、ランダムに多数の負例式をサンプル。正解を1位にしたいので、正解・負例のスコアに温度 $\tau$（例 0.07）をかけた logits に対して cross-entropy（正解インデックス 0）を最小化（InfoNCE 型）。
\end{itemize}

\subsection{ファイルと役割の対応}
\begin{itemize}[nosep]
  \item \texttt{build\_graph\_data.py}：\textbf{式--変数二部グラフ}を構築。\texttt{x} は式ノード＋変数ノードの特徴、\texttt{edge\_index} は二部エッジ、\texttt{num\_equations} と \texttt{edge\_index\_eq\_eq}（式同士の隣接）も保存。
  \item \texttt{train\_gnn\_retriever.py}：式埋め込みは GCN/GAT 出力の先頭 \texttt{num\_equations} 行を使用（二部グラフ対応）。クエリ＝SVD→MLP。\texttt{--use-gnn} / \texttt{--no-gnn} で GNN の有無、\texttt{--gnn-type gcn|gat} で GNN 種類を切替。
  \item \texttt{train\_gnn\_refine\_svd.py}：式＝SVD($E_0$) を残差 GCN で更新、クエリ＝SVD($Q$)。\texttt{edge\_index\_eq\_eq}（式同士のみ）を使用。
  \item \texttt{run\_retrieval\_experiments.py}：上記に加え \texttt{baseline\_mix}、\texttt{svd\_cos}、\texttt{q\_mlp\_to\_x}、\texttt{gnn\_refine\_svd\_residual} を同一分割で比較。二部グラフ時は式ノードのみ \texttt{x[:num\_equations]} と \texttt{edge\_index\_eq\_eq} を利用。
\end{itemize}

%----------------------------------------------------------------------
\section{今後の予定・TODO}
%----------------------------------------------------------------------
\begin{itemize}[nosep]
  \item 複数 PDF の一括処理と正規化の一括適用。
  \item 正規化ルール（\texttt{notation\_map.json}）のドメイン別拡張。
  \item 抽出結果の検証用テスト（期待する数式との差分チェック）。
  \item GitHub との同期と CI（テスト・リント）の検討。
\end{itemize}

\subsection{試したいこと・実験アイデア（大胆に）}
\label{sec:ideas}
\begin{itemize}[nosep]
  \item \textbf{GNN の種類}：GCN に加え GAT（Graph Attention）を試す。二部グラフでは「どの変数ノードを重視するか」を attention で学習できると有利な可能性。
  \item \textbf{エッジ重み}：変数ごとの IDF や、クエリ I/O 変数との一致度でエッジに重みを付け、メッセージの強さを変える。
  \item \textbf{クエリ条件付きグラフ}：推論時に「このクエリの入力・出力変数を含む式／変数」にだけメッセージを流す、あるいはエッジをマスクする。タスクに効く構造に絞る。
  \item \textbf{Hard negative}：ランダム負例に加え、変数は被るが正解でない式を意図的に負例にし、境界を学習しやすくする。
  \item \textbf{式の言い換え・拡張}：同一式の context 言い換えや変数名の正規化違いをデータ拡張し、埋め込みのロバスト性を高める。
  \item \textbf{評価の安定化}：同じ分割で複数 seed を回し、平均・標準偏差で報告。小規模データなのでばらつきを明示する。
\end{itemize}

%----------------------------------------------------------------------
\section{更新履歴}
%----------------------------------------------------------------------
\begin{itemize}[nosep]
  \item 2026-08-16（加藤先生の 7 コメント対応：進捗報告の分かりにくい箇所を修正）：
    加藤先生が \texttt{progress\_report\_2026-08-09\_SK.pdf} に付けた 7 件の注釈
    （PyMuPDF で抽出，著者 kato）に対応して報告書と較正実験を修正した．
    (1)「AUC の説明が意味不明」→ 正解式 1 本と不正解式 1 本の全組で値を比べ
    正解式が大きい組の割合と数える，という手順の形に書き直した．
    (2)「入出力だけ変えた組が 1.000 は当たり前では」→ そのとおりなので
    「説明文を変えていないので当然で，同じ説明文が 1 になる確認にとどまる」と明記し，
    内容同一の証拠は言い換えの組（0.944）であると位置づけ直した．図の行名にも（説明文は同じ）を付けた．
    (3)「無作為でなく無関係な 2 ケースと比較すべき／異なるケースで類似度が高くなるのはどんなとき？」→
    \texttt{calibrate\_similarity\_scale.py} に\textbf{無関係な組}
    （出典文献が重ならず同じ物理モデルにも由来しない組，n=191,722）を実装して基準を置き換えた
    （中央値 0.213，無作為 0.218 とほぼ同じ．JSON には両方保存）．
    高類似の無関係な組は 3.9\%（$>0.5$）で，(i) 同種の装置・現象を扱う組（CSTR どうし 0.70），
    (ii) 文献横断ケースの説明文が同じ雛形で文面がほぼ同じ組（最大 0.999）と特定し，本文に 1 項目を追加した．
    (4)(5)「意味の近さ 1 つの AUC って何？」→「意味の近さという特徴量 1 つだけで正解式を並べたときの AUC」の形に展開．
    (6)「類似度 0.437 って何？」→ 数値がすべて AUC であることを各所で明記（「文章類似度の AUC は 0.437」）．
    表 1 のキャプションに「列は不正解式の範囲の違い」を追記．
    (7)「乱数とシャッフルの関係・手順が分からない」→ \S 4 方法を手順の形に書き直した
    （対象特徴量だけ候補式間で入れ替え他は元のまま／入れ替え方を変えて 20 回平均／
    さらにデータ分割と訓練をやり直した乱数 10 通りのモデルで全体を繰り返し平均と標準偏差）．
    4 ページ維持・pytest 51 passed．
    その後の本人指示で，(2) に対応して書いた
    「入出力だけ変えた組の 1.000 は説明文を変えていないので当然」の 2 文は\textbf{不要}として削除し，
    代わりに結果の 1 項目めに\textbf{比較の基準に使う「無関係な組」の説明}
    （出典文献が重ならず，同じ物理モデルにも由来しない組）を置いた
    （方法 (a) と合わせて 2 箇所で定義し，数値のそばで基準が分かるようにした）．
  \item 2026-08-11（進捗報告の書き直し：報告内容を\textbf{2 点に絞る}）：
    \texttt{docs/progress\_report\_2026-08-09\_topic.tex} を全面的に書き直した．
    報告する内容を(1)\,意味の近さは SVD 空間で十分な区別を作れているが第 2 段では効かない，
    (2)\,PFI で各特徴量の重要度を測った，の 2 点に限定し，
    正解・不正解のケース分析（誤りの内訳・誤りの 2 型・被覆率）と
    PFI のケース別分解（分野一致度の相殺）は次回に回して本文から外した．
    実験 1 は較正の箱ひげ図（\texttt{fig\_similarity\_scale}）を主図に据え，
    候補内 AUC は表に集約した（全 DB 0.950 対 候補 50 件 0.594）．
    実験 2 は 10 特徴量と 2 群の PFI 図（\texttt{fig\_pfi\_importance}）を主図とした
    （変数の重なり 6 特徴量 $+0.678$，話題の近さ 4 特徴量 $+0.006$，$p=0.10$）．
    図の凡例・軸ラベルを報告書の用語に統一（テキスト類似→文章類似度，入出力Jaccard→変数の一致度，
    意味類似(SVD)→意味の近さ，特化度→式の特化度 等．プログラム識別子を排除）し，
    \texttt{generate\_pfi\_figure.py} と \texttt{calibrate\_similarity\_scale.py} の図サイズと表題を調整した．
    意味の近さ・式の特化度・分野の一致・分野の一致（集合版）の定義を本文に加え，自己完結させた．
    その後，本人指摘「図の文字が小さすぎる」を受けて\textbf{両図の文字を 1.5 倍}にした
    （両スクリプトに倍率定数 \texttt{FS} を導入し，目盛り・軸ラベル・凡例・値ラベル・表題の全指定に適用）．
    図の外形寸法は 1\% 以内で不変なので，報告書に同じ幅で貼ったときの実寸がそのまま 1.5 倍になる．
    PFI 図は文字拡大で(b)の縦軸ラベルが枠外へはみ出したため，
    ラベルを「Recall@K の低下」に短縮し左右のパネル幅比を 2.3:1 から 2.0:1 に変えた．
    さらに本人指示で実験 2 の節題を
    「各特徴量の重要度を PFI（Permutation Feature Importance）方法で測る」に変更し，
    略語 PFI を \S 2 の初出で導入して \S 4 本文の重複した正式名称を外した
    （英語名が節題で改行分断されないよう \texttt{\textbackslash mbox} で囲んだ）．
  \item 2026-08-09（PFI のケース別依存プロファイル：\textbf{分離度仮説の不成立}と分野一致度の\textbf{相殺}の実証）：
    順列特徴量重要度（PFI，Permutation Feature Importance）を集約値からケース単位に降ろし，
    ある特徴量を置換したとき正解から不正解へ転じるケースを同定して属性で特徴づける分析を追加した
    （\texttt{analyze\_pfi.py} を拡張し per-case 収集を統合，\texttt{analyze\_pfi\_profile.py} を新規作成）．
    設定A・10 シード・20 置換で per-case 50,534 記録を出力．
    集約 \texttt{pfi\_results.json} は既存版と\textbf{バイト単位で同一}であり，集約 PFI と乱数系列の保持を実証した．
    \textbf{当初仮説（分離度 sep が大きいケースほど反転しやすい）は成立しなかった}．
    sep（正解式と不正解式の特徴値の差）と期待 Recall 低下の Spearman 順位相関は
    補完性 $\rho=-0.011$（$p=0.78$，非有意），一貫性 $\rho=+0.105$（$p=0.005$），
    分野一致度 $\rho=+0.243$（$p=2.9\times10^{-11}$，$n=727$ ケース）で，
    \textbf{効かないとしてきた分野一致度が最も強い単調性を示した}．
    補完性の依存ケースは sep がむしろ低い（標準化平均差 $-0.31$）．
    代わりに頑健に現れたのは\textbf{系の大きさ}で，依存ケースは正解式数・入力変数数・出力変数数・出典数が
    いずれも大きい（一貫性の正解式数で標準化平均差 $+0.49$）．
    分野一致度の集約 PFI が $\approx 0$ なのは寄与が無いためではなく\textbf{相殺}による：
    全ケース母集団で分離時（sep $>0$）の平均低下 $+0.012$ に対し逆分離時（sep $<0$）は $-0.034$ と符号が反転する
    （Mann--Whitney $p=1.3\times10^{-8}$）．すなわち置換は分離時には性能を下げ，逆分離時には上げる．
    なお解析途上で\textbf{疑似反復}（同一ケースが複数シードのテストに出現し，2,067 記録が実は 727 ケース，平均 2.84 回）を
    独立観測として検定していた誤りを発見し，ケース単位に集約してから検定するよう修正した．
    修正前は補完性が $p=0.009$ と有意に見えていたが，修正後は $p=0.78$ で有意性は消えた．
    また依存群と頑健群に同一ケースが重複していた（補完性で 62.5\,\%）問題も同時に解消した．
    結果：\texttt{experiments/pfi\_per\_case.json}，\texttt{experiments/pfi\_profile\_stats.json}，
    \texttt{docs/figures/fig\_pfi\_dependence\_profile.png}．
  \item 2026-06-03（Stage2 reranker 強化の試行と\textbf{負の結果}）：
    厳格 Recall@K\_correct の向上を狙い，\texttt{set\_aware\_reranker.py} に
    (1) hard negative マイニング（\texttt{-{}-hard-neg}：現在スコア上位の negative を採用），
    (2) listwise 損失（\texttt{-{}-loss listwise}：multi-positive softmax CE），
    (3) 出力補完特徴 gOutNov（候補が参照集合に無い出力変数を固有に覆う割合）と
    新モード \texttt{reranker-10S2}（gDom$\to$gOutNov）・\texttt{reranker-11S}（set-aware 4 特徴量）を追加した（既存 reranker-10S は不変・後方互換）．
    DAE・3 シードで比較した結果，\textbf{全構成が Recall@K\_correct $=0.771$--$0.776$（差 $\le 0.006$）でシード標準偏差 $\pm 0.07$ の内側}に収まり，
    hard negative・listwise・gOutNov のいずれも有意な改善を生まなかった．Recall@20 は $0.91$ で候補自体は上位に揃うため，
    律速は候補を\textbf{独立にスコアする}現行 MLP が\textbf{集合としての選択}を表現できない点にあると診断．
    次の本質的改善は候補集合の相互作用を陽に扱う\textbf{構造的拡張}（候補集合への注意機構／GNN，逐次的集合構成）と結論した．
    結果：\texttt{experiments/stage2\_\{feat,hardneg,listwise\}.json}．
    \textbf{（追記：greedy 逐次選択は成功）}上記の負の結果を受け，診断「集合選択が律速」を直接検証するため
    \texttt{set\_aware\_reranker.py} に\textbf{推論時 greedy 集合構築}（\texttt{-{}-greedy}：1 式選ぶごとに
    set-aware 特徴量を既選択集合を参照に再計算して次を選ぶ）と関数 \texttt{greedy\_order}，
    \texttt{compute\_set\_aware\_features} の \texttt{ref\_indices} 引数を追加した（学習済みモデルをそのまま使用，新規学習不要）．
    DAE・3 シードで \textbf{reranker-10S/10S2/11S すべてが $+0.015$--$0.018$ 一貫して改善}（\texttt{stage2\_greedy.json}）．
    効いた 2 レバー（top-$k$ 拡大と greedy）を結合し，reranker-10S・5 シードで
    \textbf{Recall@K\_correct $0.745\to0.792$（$+0.048$，$+6.4\%$），Recall@20 $0.887\to0.953$}を達成
    （\texttt{cap\_ref.json}, \texttt{cap\_greedy200.json}）．損失・負例の改良では届かなかった厳格指標を「集合としての選択」が動かすことを実証．
    図 H（\texttt{docs/figures/fig\_v6\_stage2.png}）を追加し v6 §精度向上を更新．次は greedy の学習段階導入（teacher forcing）が課題．
    \textbf{（追記2：near-miss 分析）}「漏れた正解がどれだけ上位に詰まっているか」を測る指標を新設．
    \texttt{evaluate\_multi\_eq.py} の \texttt{case\_metrics} が各正解式の順位 \texttt{ranks} を返すようにし，
    \texttt{generate\_figures\_v6.py} の \texttt{fig\_nearmiss} で\textbf{回復曲線 Recall@$(K+m)$}（正解数 $K$ に余裕 $m$ を足した Recall，$m{=}0$ が Recall@K\_correct）と
    \textbf{超過順位（順位 $-K$）の累積分布}を作図（候補 200 件予算で統一）．DAE・top-200・3 シードで baseline／reranker-10S／+greedy を比較した結果，
    \textbf{reranker・greedy は漏れた正解の 45--47\% を上位 3 件以内に詰める}（baseline 18\%），$K+3$ で Recall $\approx 0.9$．
    本手法の誤りの大半が near-miss であり，残る far-miss（漏れの 12--13\% が候補 200 件外）が Stage1 強化の余地と判明．
    図 I（\texttt{fig\_v6\_nearmiss.png}）を v6 §精度向上に追加．データ：\texttt{experiments/ranks\_std.json}, \texttt{experiments/ranks\_greedy.json}．
    \textbf{（追記3：ゼミ報告 2026-06-14 作成）}前回ゼミ（2026-05-19，3{,}244式・1{,}338ケース・MRR 0.963）からの進捗をまとめた新ゼミ報告
    \texttt{docs/progress\_report\_2026-06-14\_seminar.tex/.pdf}（7 ページ，前回と同形式，日付 2026-06-14）を作成．
    構成：前回レビュー $\to$ データ大規模化（11{,}146式・2{,}838ケース）$\to$ 評価指標精緻化（Recall@K\_correct）$\to$
    評価結果（難易度勾配の 3 者比較・文献横断 LLM 比較・議論点(1)への回答=仮説B・頑健性）$\to$ 精度向上（greedy・near-miss）$\to$ 今後．
    前回の議論点(1)「関連性の低い式も精度向上に機能するか」へは，フル DB vs 限定 DB（$+0.048$，$p=0.044$）で\textbf{回答}した．
    Ultracode の検証 workflow（9 エージェント並列）で全数値を実験 JSON に敵対的照合し \textbf{86 件 MATCH}，指摘 23 件
    （用語の和文説明・文献横断 LLM 比較の補完・Recall@20 の表記統一・記号 K と top-$k$ の区別等）を修正した．

  \item 2026-06-02（進捗報告 v6 作成：データ拡大・難易度勾配・仮説B再検定・頑健性・LLM 同一条件比較の統合）：
    前回 v5（3{,}244式・1{,}107ケース）以降の Phase 3--5 成果と外部比較を\textbf{差分版の進捗報告 v6}としてまとめ，
    \texttt{docs/progress\_report\_2026-06-02\_v6.tex}（lualatex / ltjsarticle，8ページ）を新規作成．
    主な内容：(1) データ規模 11{,}146式・2{,}838ケース（DAE 1{,}000件含む），
    (2) 難易度勾配で本手法の Recall@K\_correct 優位が標準 $+0.127$ $\to$ +DAE $+0.210$ $\to$ DAEのみ $+0.217$ と拡大，
    (3) 仮説B再検定でフル DB が Recall@K\_correct で有意（$+0.048$，$p=0.044$，Cohen's $d=0.74$ 中；MRR\_first では飽和し差なし）と判明し v5 の暫定結論「DB 規模の影響なし」を訂正，
    (4) ハイパラ sweep 13 設定で Recall@K\_correct レンジ $0.0094$（シード間 std の約 1/10）と頑健性を確認，
    (5) \textbf{外部ベースライン（LLM 直接生成）との同一条件比較を現行 11{,}146式 DB で整合}．
    \textbf{LLM 比較の再評価手順}：LLM 評価済みの単一ソース 29 件・複数ソース 20 件サブセットに対し，
    \texttt{set\_aware\_reranker.py -{}-modes baseline,reranker-10S -{}-seeds 10 -{}-save-per-case}
    で自然 5 variant（original / context\_paraphrased / random\_io\_from\_models / multisource\_original / multisource\_random\_io）学習・per-case 出力を生成し（\texttt{experiments/phase6\_full\_percase.json}），
    各対象ケースを\textbf{テスト分割に入ったシードのみ}で集計（リーク防止の held-out 評価）する \texttt{aggregate\_llm\_subset\_v6.py} を新規作成（出力 \texttt{experiments/phase6\_llm\_subset\_comparison.json}）．
    結果（Recall@K\_correct）：単一ソースは baseline $0.982$／reranker-10S $0.975$ がいずれも LLM $0.707$ を圧倒，
    複数ソースは baseline $0.575 <$ LLM $0.842 <$ \textbf{reranker-10S $0.942$} となり，
    古典 IR 単独では複数文献の組合せに不足し Set-aware リランカーが本質的に必要であることを同一条件で実証．
    なお和文太字は明朝に太字ウェイトが無いため \texttt{BoldFont} を角ゴシックに設定した．
    \textbf{（同日追記）}v5 の構成を参考に v6 を\textbf{図駆動で簡潔化}（8 $\to$ 5 ページ；目次・重複表を図に統合）し，
    機械学習研究者視点の図 4 点を新規作成（\texttt{generate\_figures\_v6.py} $\to$ \texttt{docs/figures/fig\_v6\_\{dataset,gradient,impact,llm\}.png}）：
    (A) データセット内訳・難易度分布，(B) 難易度勾配〔中核〕，(C) 正解式数別の改善＋指標飽和（仮説B）の対比，(D) LLM 三者比較．
    各図はシード分散のエラーバー・差分注記・統一配色（baseline 灰／本手法 青／LLM 橙）で「1 図 1 主張」とした．
    \textbf{（同日追記2：LLM 難ケース比較・精度向上・アブレーション）}
    加藤先生の「LLM で解けるなら不要では」「図B に LLM も並べよ」とのご指摘を受け，\texttt{evaluate\_llm\_dae.py} を新規作成し
    DAE 難ケースで LLM 直接生成（Claude）を評価．\texttt{generate\_figures\_v6.py} の図B を\textbf{3 者比較}（baseline／LLM／reranker-10S）に更新し，
    \textbf{全難易度で reranker $>$ baseline $>$ LLM}（最難 DAE で LLM $0.264$ に対し本手法 $0.699$，約 2.6 倍）を示した．
    精度向上として Stage1 の retrieval 天井を診断（DAE Recall@50 $=0.86$，高 X で取りこぼし $X{=}10$ で $0.68$）し，
    候補数 top-$k$ を $50\to200$ に拡大して reranker-10S を再評価，\textbf{実用 Recall@20 を $0.912\to0.968$ に改善}した
    （\texttt{topk\_dae\_\{50,100,200\}.json}；厳格 Recall@K\_correct は $+0.02$ で，Stage2 の順位付けが律速と判明）．
    Set-aware 3 特徴量のアブレーション（\texttt{ablation\_setaware.json}）で \textbf{gComp・gCoh が主力，gDom は寄与薄}と判明．
    図 E（シード分散）・F（アブレーション）・G（精度向上）を追加し v6 は 7 ページに．画像内タイトルの図番号は LaTeX キャプション採番へ一本化した．

  \item 2026-06-01（Phase 5 完了：reranker ハイパーパラメータ sweep と頑健性の確認）：
    \texttt{set\_aware\_reranker.py} に \texttt{-{}-hidden-dim}, \texttt{-{}-margin}, \texttt{-{}-batch-size}, \texttt{-{}-n-neg-samples}, \texttt{-{}-weight-decay} を追加し，
    sweep 実行用 \texttt{run\_phase5\_sweep.py}（config を子プロセスで並列実行・結果集約）と
    分析用 \texttt{analyze\_phase5\_sweep.py}（ランキング・パラメータ感度・best 抽出，複数 sweep を \texttt{-{}-config-glob} で横断結合可）を新規作成．
    GPU サーバ初期化補助 \texttt{setup\_gpu\_server.sh}（rsync + Docker コンテナ + 依存導入 + CUDA 確認 + sweep キックオフ）も作成．
    \textbf{実行構成}：固定軸 \texttt{reranker-10S} + 10 seed + 訓練データ（original + multisource\_v3 + dae\_X*），
    探索軸 lr $\in \{$1e-4, 3e-4, 1e-3, 3e-3, 1e-2$\}$，hidden\_dim $\in \{$64, 128, 256$\}$，epochs $\in \{$20, 40$\}$，margin $\in \{$0.1, 0.3$\}$．
    計 13 config を評価（A100$\times$4 サーバ 8 config + ローカル lr sweep 5 config）．
    \textbf{結果（主指標 Recall@K\_correct）}：全 13 config が $0.7321$--$0.7415$ に収まり，\textbf{config 間レンジ $0.0094$ は seed 間標準偏差 $0.1008$ の約 1/10}．
    すなわち探索したいずれのハイパーパラメータ（lr を 100 倍振っても，hidden\_dim・epochs・margin を変えても）\textbf{統計的に有意な差を生まない}．
    best config は lr$=$3e-4, hidden\_dim$=$64, epochs$=$40, margin$=$0.1（R@K\_c $=0.7415$）だが既定値（lr$=$1e-3, hidden$=$64, epochs$=$15）との差は誤差範囲．
    \textbf{結論}：reranker-10S の強さはハイパーパラメータの微調整に依存せず\textbf{頑健}であり，「脆弱なチューニング由来の見かけ上の性能」ではないことを定量的に確認した（査読対策として重要）．
    \textbf{運用上の知見}：本手法は MLP が小さく \textbf{GPU 加速の恩恵がほぼ無い CPU-bound} 処理であり，共有 GPU サーバ（過負荷時 load 53/16 コア）より\textbf{非競合のローカル実行の方が予測可能}．長時間 CPU ジョブは \texttt{caffeinate -i} でスリープ抑止が必須．
    保存先：\texttt{experiments/phase5\_sweep/}（config 個別 JSON），\texttt{experiments/phase5\_best.json}．

  \item 2026-05-29（Phase 4 完了：特性別性能分析 + 仮説 B 再検定）：
    \textbf{(1) 仮説 B（Full DB vs Narrow DB の本手法 MRR 差）を 3 指標で再検定．}
    旧結果は \texttt{MRR\_first} のみで $p=0.78$（n.s.）と「DB 規模の影響なし」と結論したが，
    \texttt{MRR\_first} は天井 1.0 付近で飽和する指標であった．新指標 \texttt{Recall@K\_correct}（$K=$正解式数 $n_{\text{correct}}$，R-Precision と同義）で同検定を行うと，
    \textbf{Full DB} $0.9229$ vs \textbf{Narrow DB} $0.8749$，差 $+0.0480$（paired t: $t=2.34$, $p=0.044$；Wilcoxon $p=0.037$；Cohen's $d=0.74$ medium）と
    \textbf{中程度効果の有意差}を確認．\textbf{MAP} も $0.953 \to 0.931$（$p=0.096$，$d=0.59$ medium）と境界域で同方向の差．これにより
    「サチった指標が差を隠していた」ことが明確になり，DB 拡張の効果は実際にはあるが評価指標の選択によって見え方が変わる，という重要な知見を得た．保存先：\texttt{experiments/phase4\_hypothesis\_b\_test.json}．
    \textbf{(2) 特性別 (stratified) 性能分析．}
    \texttt{set\_aware\_reranker.py} に \texttt{-{}-save-per-case} フラグを追加して各ケースの評価結果を JSON にダンプする機能を実装し，
    10 seed・1,592 ケース・1,223 ユニーク（original + multisource\_v3 + dae\_X1..X10）について
    \texttt{baseline} vs \texttt{reranker-10S} を 4 軸（正解式数 $n_{\text{correct}}$，横断ソース数 $n_{\text{sources}}$，
    入力変数数 $n_{\text{input}}$，出力変数数 $n_{\text{output}}$）で層別化．
    主要結果（差は \texttt{Recall@K\_correct}）：
    $n_{\text{correct}}=1$ では $+0.02$ と差なしだが，$n_{\text{correct}}=2$ で $+0.21$，$=3$ で $+0.26$，$5\text{-}10$ で $+0.24$ と
    \textbf{複数式ケースで大幅改善}．$n_{\text{input}}=1\text{-}3$ では $-0.02$（baseline がすでに $0.92$）で本手法不要，$11+$ で $+0.25$（baseline $0.45 \to$ reranker $0.70$）．
    \textbf{Recall@K\_correct} は MRR\_first の天井問題を解消し，提案手法の真の改善幅（baseline $0.526 \to$ reranker $0.735$，$+0.21$）を露わにした．
    \texttt{analyze\_phase4\_characteristic.py} と \texttt{analyze\_phase4.py} を新規作成・更新．保存先：\texttt{experiments/phase4\_characteristic\_analysis.json}，図 \texttt{docs/figures/fig\_phase4\_stratified.png}．

  \item 2026-05-28（Phase 3 完了：制約付き実験 A/B + DAE データセット）：
    \textbf{(1) DAE 訓練ケース 1,000 件追加}．\texttt{generate\_dae\_cases.py} を新規作成し，Anthropic API（Claude Sonnet 4.5）で
    $X \in \{1, 2, \ldots, 10\}$ の各値について 100 件ずつ生成（合計 1,000 件，\texttt{variant\_type = dae\_X\{N\}}）．
    アルゴリズム：(i) ランダムに ODE を 1 本選択し正解集合の初期種とする，(ii) その微分変数を出力初期値に追加，(iii) 同一ソース内の変数共有式を優先的に正解集合に追加，(iv) $|$正解集合$|=X$ になるまで反復，(v) $X>1$ の場合は追加された各式の代表変数を出力に加えて自由度 0 を保証，(vi) 残り変数を入力とし context を Claude API で生成．
    最終訓練ケース数：$1{,}838 \to 2{,}838$．
    \textbf{(2) Phase 3 制約付き実験}．\texttt{set\_aware\_reranker.py} に \texttt{-{}-variants} フィルタ（プレフィックスマッチ対応）と \texttt{-{}-output} オプションを追加．
    実験 A（\texttt{original + multisource\_v3 + dae\_X*}）と B（\texttt{dae\_X*} のみ）を 10 seed で実行．
    結果：A で baseline MRR\_first $0.8024 \to$ reranker $0.9232$（$+0.121$），B で baseline $0.7806 \to$ reranker $0.9057$（$+0.125$）．
    \textbf{DAE のみ}に絞ると \textbf{ベースラインの MRR\_first が $0.78$ まで低下}（フルデータでは $0.917$）し，本手法の改善幅が拡大（MAP $+0.30$）．
    DAE 系ケースが従来 random\_io 拡張より格段に \textbf{難しいデータ}であることが定量的に示された．
    \textbf{(3) Phase 1 評価指標拡張}．\texttt{evaluate\_multi\_eq.py} に \texttt{Recall@K\_correct}（$K=n_{\text{correct}}$）を追加．従来の \texttt{MRR\_first} は飽和しやすく改善幅が見えにくい問題を解消．
    Phase 3 結果での Recall@K\_correct：A で $0.525 \to 0.735$（$+0.210$，\textbf{+40\% 相対}），B で $0.481 \to 0.699$（$+0.218$，\textbf{+45\% 相対}）．
    関連スクリプト・データ：\texttt{generate\_dae\_cases.py}, \texttt{run\_phase3\_experiments.sh}, \texttt{experiments/phase3\_A\_original\_multi\_dae.json}, \texttt{experiments/phase3\_B\_dae\_only.json}．

  \item 2026-04-18（自由度検証済1,107ケースでの主要結果の再評価）：\texttt{two\_stage\_query\_conditioned.py}を5 seedで再実行．\textbf{reranker-10の優位性が消失}し，reranker-7が最良となる大きな結果の変化を観測（baseline $0.924\to0.937$，reranker-7 $0.937\to\mathbf{0.959}$，reranker-10 $0.949\to0.939$）．全体として各手法のMRRが向上したことから，自由度検証により学習・評価双方の品質が改善していることが示された．reranker-10のグラフ特徴量（\#7--9）は初回データの不正ケース（物理的に解けないケース）で正解順位を押し上げる役割を果たしていた可能性があり，現行の特徴量設計の見直しが必要．進捗レポートv2§5に新しい結果表・バリアント別・初回データとの差分を追加，§6考察に再解釈を追記，§7今後の課題に「グラフ特徴量の再設計」を高優先度で記載．実験結果は\texttt{experiments/query\_conditioned\_comparison.json}に保存．

  \item 2026-04-19（追補レポート作成：自由度解析 + Set-aware 特徴量）：前回報告（v2）後の一連の研究成果を追補報告としてまとめ，\texttt{docs/progress\_report\_2026-04-19\_addendum.tex} を新規作成．主な内容：(1) 自由度解析による訓練データ品質改善（1{,}107件を全件可解に再構成），(2) 既存グラフ特徴量 f7--f9 の診断（f9 は AUC 0.087 で逆信号など）と限界の明確化，(3) Set-aware 3特徴量（gComp=補完性，gCoh=一貫性，gDom=ドメイン一致）の設計と評価，(4) 複数式評価指標（MRR\_worst, Recall@K, MAP）の導入．10 seed paired t-test の結果：MAP $0.936 \to 0.955$（$p=0.004$），複数式 Recall@3（旧 FullRecall@3）$0.628 \to 0.693$（$p=0.032$），複数式 MRR\_worst $0.376 \to 0.397$（$p=0.024$）．reranker-10S（基本7 + set-aware 3）が統計的に有意な改善を達成．論文化の方向性を「Set-aware graph features for multi-equation retrieval」に確定．関連スクリプト：\texttt{filter\_solvable\_cases.py}, \texttt{expand\_solvable\_random\_io.py}, \texttt{diagnose\_graph\_features.py}, \texttt{graph\_feature\_ablation.py}, \texttt{evaluate\_multi\_eq.py}, \texttt{set\_aware\_reranker.py}．実験結果：\texttt{experiments/graph\_feature\_diagnosis.json}, \texttt{experiments/graph\_ablation\_results.json}, \texttt{experiments/multi\_eq\_evaluation.json}, \texttt{experiments/set\_aware\_results.json}．\textbf{2026-05-08 追記}：FullRecall@K（全正解上位K内割合）から標準的な Recall@K（$|R_q \cap \mathrm{Top}_K|/|R_q|$ の平均）に指標を切替．\texttt{evaluate\_multi\_eq.py} に Recall@K 計算を追加し再実行．

  \item 2026-04-18（自由度=0を保証するrandom\_ioの追加生成）：自由度フィルタで707件に減少した訓練データを，\textbf{自由度=0を構成上保証する新規random\_io}を追加生成することで1,107件に復元．\texttt{expand\_solvable\_random\_io.py} を新規作成し，各コアの正解式変数集合$V$と方程式数$n_{\text{eq}}$に対し，$V$から$n_{\text{eq}}$個を出力として選び残りを入力とする全組合せから，既存の入出力組と重複しないものを採用する方式で400件生成．新規400件全てが自由度解析をパス．最終内訳：original 93, paraphrased 279, swap\_io 2, random\_io 733（元333＋新規400）＝1,107件．進捗レポートv2§3.4.2を「自由度解析＋可解ケース追加生成」に再改訂，§2.1概要表・§4.5正例・§5結果注記・§7今後の課題の数値を1,107に更新．

  \item 2026-04-18（自由度解析による訓練ケースのフィルタ実装）：加藤先生指摘「物理的に解けないケースは学習に使ってはいけない」への対応．\texttt{filter\_solvable\_cases.py} を新規作成し，各訓練ケースについて正解式集合の変数数と独立方程式数を照合する自由度解析を実装．判定条件：(i) 入出力変数が全て式集合内にあること，(ii) 出力変数が未知側にあること，(iii) $|\text{未知}|\le$方程式数（relaxed：過剰決定系も物理的に解けるため許容）．訓練データ1{,}107ケースに適用した結果，707ケース（63.9\%）のみが可解と判定されたため，\texttt{training\_cases.json} を差し替え（元データは \texttt{training\_cases.pre\_dof\_filter.bak.json} に退避）．内訳：original 123$\to$93, context\_paraphrased 369$\to$279, swap\_io 123$\to$2（98\%除外），random\_io 492$\to$333．swap\_ioは入出力非対称の元ケースが大多数であったため激減．進捗レポートv2§3.4.2を「自由度解析の実装済み」へ改訂，§2.1概要表・§4.5正例定義・§5結果の注記・§7今後の課題に反映．MRR再評価は今後の課題として明示．

  \item 2026-04-17（進捗レポートv2の§2.1拡充：変数表記揺れの方針明記）：前回加藤先生に「変数表記の正規化に取り組む」と述べた件について，実データ（3{,}244式）を精査した結果を受け，\textbf{現段階では全面正規化を行わない判断}に改めた旨を\texttt{progress\_report\_2026-04-13\_v2.tex}の§2.1に明文化．表記揺れの実態（比熱が$c_p$/$C_p$/$c_w$/$Cp$の4種99回、前指数因子が3種31回、密度が3種251回）を表で示し，全面正規化を行わない3つの根拠（(1)記号の多義性による誤正規化リスク、(2)文献単位の分割評価下では同文献内の表記が一貫し効果限定的、(3)現MRR 0.949で誤差の主因でない可能性）を明記．実装済みの\texttt{normalize\_notation.py}（3物理量ルール）によるアブレーション実験を今後の内部検証として位置づけ．レポート本文では具体的なファイル名への言及は削除し，長期方針（各変数の物理的意味の自然言語記述を埋め込み空間で比較し，意味的に近いノード同士をグラフ上でソフトに接続する学習的アプローチ）のみを記載．

  \item 2026-04-17（進捗レポートv2への注釈付与機構の刷新）：v2 PDFへのハイライト＋スティッキーノート付与を，従来の「固定文字列マッチ」方式から\textbf{提出版PDFとのテキスト単語レベル差分}方式に全面変更（\texttt{docs/annotate\_report\_v2.py}を新規作成，旧\texttt{add\_annotations.py}/\texttt{add\_annotations2.py}を置換）．\texttt{pymupdf}の\texttt{get\_text("words")}で両PDFの単語列を取り，\texttt{difflib.SequenceMatcher}で差分を取ることで，\textbf{実際に変更された箇所のみを正確にハイライト}（タイトルなど未変更箇所の誤ハイライトを解消）．正規化は Unicode NFC、ダッシュ類・引用符類・全角句読点・各種空白を統一し表記揺れを吸収．結果：321箇所のハイライト（equal=1{,}234単語・insert=217単語・replace=242単語）と25箇所のスティッキーノート（各修正に対応する加藤先生コメント内容を詳述）．

  \item 2026-04-09（Phase 1 Round 2: 3{,}244式への拡大 + 再評価）：\texttt{expand\_dataset.py} にハンドブック第2弾（新10ドメイン＋既存6ドメインの応用式、計320式）と大型 PDF 分割抽出（\texttt{pypdf} で80ページずつ分割、3冊39チャンク）を追加。結果：\textbf{1{,}613式 $\to$ 3{,}244式}、ソース $52 \to 71$、変数 $2{,}098 \to 3{,}704$。\texttt{generate\_training\_cases.py} を大規模カタログ対応に改修（ドメイン別バッチ分割＋503リトライ）。123コア $\to$ 1{,}107訓練ケースで5\,seed 比較：
  \begin{itemize}[nosep]
    \item baseline: MRR $0.924 \pm 0.040$, original $0.921$
    \item reranker-7: MRR $0.937 \pm 0.018$, original $0.942$
    \item reranker-10 (graph): MRR $\mathbf{0.949 \pm 0.031}$, original $\mathbf{0.952}$
  \end{itemize}
  \textbf{核心的発見}：(a) グラフ特徴量（reranker-10）が全指標・全 variant で最良、(b) reranker-7 を明確に上回る（$+1.2\%$ overall, $+1.0\%$ original）、(c) 975式$\to$3{,}244式のスケーリングでグラフ特徴量の有効性が単調増加。「データ規模がグラフ構造の価値を解放する」仮説を強く支持。Gemini API キーをハードコードから環境変数専用に変更済み。

  \item 2026-03-31（Phase 1: データ拡大 + 再評価）：\texttt{expand\_dataset.py} を新規作成し、3つの手法でデータセットを拡大。(1)\,ハンドブック式の大幅拡張：LLM 知識から15新ドメイン（流体力学・伝熱・物質移動・熱力学・反応工学・プロセス制御・電気化学・バイオプロセス・環境工学・振動・電気工学・高分子・燃焼・分離・輸送現象）で330式を生成。(2)\,Semantic Scholar API で新論文を探索、open-access PDF 17本をダウンロード。(3)\,新規+既存未処理 PDF 15本から308式を抽出。結果：\textbf{975式 $\to$ 1{,}613式}（$+65\%$）、ソース $22 \to 52$。訓練ケースも再生成（$154 \to 187$ コア、$1{,}343 \to 1{,}680$ 拡張）。

  拡大データでの再評価結果（5\,seed）：
  \begin{itemize}[nosep]
    \item baseline: MRR $0.900 \pm 0.069$, original $0.885$（旧: $0.875, 0.839$）
    \item reranker-7: MRR $0.942 \pm 0.052$, original $0.936$（旧: $0.885, 0.860$）
    \item reranker-10 (graph): MRR $\mathbf{0.948 \pm 0.046}$, original $0.935$（旧: $0.879, 0.842$）
  \end{itemize}
  \textbf{核心的発見}：データ拡大により (a) リランカーがベースラインを明確に上回る（$+4\text{--}5\%$、旧 $+1\%$ 未満）、(b) グラフ特徴量が初めて悪化せず全体 MRR で最良、(c) 分散が半減（$0.119 \to 0.052$）。「モデル改善より先にデータの質と量が重要」という仮説を実証。また、ハードコードされた Gemini API キーが GitHub 公開でリーク報告されたため、全5ファイルから削除し環境変数のみに変更。

  \item 2026-03-30（Phase 3b: クエリ条件付きグラフ特徴量）：\texttt{two\_stage\_query\_conditioned.py} を新規作成。二部グラフ（式--変数）の構造情報をクエリ条件付きで特徴化する3指標を設計：(1)\,2-hop reachability（クエリ変数から候補式への到達率）、(2)\,graph neighbor overlap（候補式のグラフ近傍とクエリ変数関連式の Jaccard）、(3)\,indirect bridge ratio（候補変数がクエリ変数関連式にブリッジする比率）。7次元リランカーに3特徴を追加した10次元リランカーで5\,seed 比較。結果：\textbf{クエリ条件付き特徴量も改善に寄与せず}。MRR $0.879 \pm 0.123$（7次元の $0.883 \pm 0.119$ から微低下）、original MRR $0.842$（7次元の $0.860$ から $-0.018$）。GCN ベース特徴量（Phase 3）と同様のパターン。\textbf{結論}：現データ規模（975式・154 original ケース）ではグラフ構造が有効な追加信号を提供しない。7次元学習リランカー（no-gnn）が最良手法として確定。論文では「なぜグラフが効かないか」の分析を contribution に含める。
  \item 2026-03-30（Phase 3: GNN 特徴量統合）：\texttt{two\_stage\_retrieval\_gnn.py} を新規作成し、3方式を5\,seed で比較。(1)\,baseline\_mix（固定重み）、(2)\,no-gnn リランカー（7次元特徴、学習済み MLP）、(3)\,gnn リランカー（+GCN-refined SVD/CodeBERT 特徴の9次元）。結果：\textbf{no-gnn リランカーが最良}（MRR $0.885 \pm 0.113$、original $0.857$）。GNN 特徴を追加すると分散が増大し性能低下（MRR $0.873 \pm 0.148$）。原因分析：(a) CodeBERT GCN のセルフシフト特徴はケース依存情報を持たない、(b) SVD GCN の学習が不十分（5 epoch, 64 sample）、(c) グラフ構造がクエリ固有の情報を反映していない。\textbf{結論}：特徴の重み付けを固定→学習に変えるだけで改善が得られるが、現状のグラフ構造から有効な特徴を抽出するにはクエリ条件付き GNN が必要。
  \item 2026-03-30（Phase 2: 2段階 retrieval）：\texttt{two\_stage\_retrieval.py} を新規作成。Stage\,1 で \texttt{baseline\_mix} が Top-$K$（$K\!=\!50$）候補を高速に取得し、Stage\,2 でニューラルリランカー（7次元特徴量、3層 MLP、ペアワイズ margin loss）が精密にリランキング。5\,seed 評価の結果、全体 MRR $0.875 \to 0.881$（$+0.006$）。\textbf{original ケースでは} MRR $0.839 \to 0.856$（$+0.017$）と改善を確認。ただし seed 456 で劣化するなど分散が大きく、統計的有意性は未達。
  \item 2026-03-30（データ品質修正）：\texttt{fix\_dataset\_quality.py} を新規作成。暗黙変数（$s$, $t$）を277式に追加、重複ケース43件を除去（$1386 \to 1343$）。バックアップを \texttt{*.bak.json} に保存。
  \item 2026-03-30（Phase 0: 評価基盤の確立）：論文化を見据え、研究の土台を再設計。
    \begin{itemize}[nosep]
      \item \textbf{ベースライン失敗分析}（\texttt{analyze\_baseline\_failures.py} 新規作成）：全1386ケースで \texttt{baseline\_mix} を評価し、失敗ケース（RR $<0.5$）の85件（6.1\%）を4パターンに分類。\textbf{cross\_domain}（100\%）と\textbf{context\_divergence}（78\%）が支配的 --- クエリと式の語彙が異なるケースでベースラインが失敗しており、意味理解（GNN/ニューラル）が有効な領域を特定。可視化を \texttt{figures/failure\_analysis\_summary.png} に出力。
      \item \textbf{複数seed評価の導入}（\texttt{run\_retrieval\_experiments.py} 改修）：5 seed（42, 123, 456, 789, 1024）で source\_id 分割を変え、平均$\pm$標準偏差と95\% bootstrap CIを報告する \texttt{--seeds N} モードを追加。結果：\texttt{baseline\_mix} の Test MRR $= 0.890 \pm 0.095$（95\% CI $[0.804, 0.949]$）。単一 seed（0.899）では見えなかった高分散を確認。
      \item \textbf{variant\_type 別評価の標準化}：original ケースの MRR $= 0.851 \pm 0.080$ に対し、random\_io\_from\_models の MRR $= 0.925 \pm 0.112$ と合成ケースが有意に簡単であることを定量的に確認。論文では original ケースの指標を主要報告とする方針。
    \end{itemize}
  \item 2026-03-30（TeX）：\texttt{development\_log.tex} の \texttt{inputenc} を UTF-8 のみに整理し、更新履歴の Markdown 風バッククォートをやめてファイル名をタイプライタ体コマンドで囲む表記に修正、\texttt{uplatex} でコンパイル可能に。TeX は \texttt{brew install texlive} で \texttt{pdflatex}/\texttt{uplatex}/\texttt{dvipdfmx} を利用（手順は \texttt{docs/README.md}）。
  \item 2026-03-30：Claude Code 向け \texttt{CLAUDE.md} と \texttt{research\_context.md} のポインタを追加。\nameref{sec:file-roles}を新設し各ファイルの役割を記載。技術メモの \texttt{build\_graph\_data.py} を二部グラフ版の説明に修正し当該節へ参照。
  \item 2026-03-18（試したいこと・GAT 追加）：開発記録に「試したいこと・実験アイデア」を追加（GAT、エッジ重み、クエリ条件付きグラフ、hard negative、データ拡張、複数 seed 評価）。\texttt{train\_gnn\_retriever.py} に GAT エンコーダと \texttt{--gnn-type gcn|gat} を追加。
  \item 2026-03-18（グラフ構築の改善）：現状の「式＝ノード・共通変数＝エッジ」では GNN の強みを活かせないという分析を開発記録に全文追記。\textbf{式--変数二部グラフ}に変更（\texttt{build\_graph\_data.py}）。式ノードと変数ノードを明示し、エッジは「式が変数を含む」のみ。式埋め込みは GCN 出力の先頭 \texttt{num\_equations} 行を使用。\texttt{edge\_index\_eq\_eq} で後方互換を確保。\texttt{train\_gnn\_retriever.py}、\texttt{run\_retrieval\_experiments.py}、\texttt{train\_gnn\_refine\_svd.py}、\texttt{analyze\_dataset.py} を二部グラフ対応に修正。
  \item 2026-03-18（GNN・埋め込みの説明）：開発記録に「GNNの仕組みと埋め込み」セクションを追加。グラフ構築（ノード＝式・エッジ＝変数共有・特徴＝CodeBERT）、式埋め込みの2系統（CodeBERT \texttt{x} と TF-IDF+SVD）、GNN の2パターン（\texttt{x} を GCN で更新する方式と SVD 空間での残差 GCN）、スコアリング・学習の流れ、各スクリプトの役割を記載。
  \item 2026-03-18（方式比較とハイパラ探索）：retrieval方式（\texttt{baseline\_mix}, \texttt{svd\_cos}, \texttt{q\_mlp\_to\_x}, \texttt{gnn\_refine\_svd\_residual}）を同一の train/val/test（source\_id 単位）で比較し、学習率・重み減衰のグリッド探索を自動実行するランナー \texttt{run\_retrieval\_experiments.py} を追加。結果は \texttt{experiments/retrieval\_experiments.md} に表で出力し、現状では \texttt{baseline\_mix}（TF-IDF+変数Jaccard）が最良であることを確認。
  \item 2026-03-18（GNN次ステップ）：SVD+I/O 文字列ベースラインを維持したままグラフ構造の寄与を検証するため、SVD空間上で式埋め込みをGCNで\textbf{残差的に}更新する実験スクリプト \texttt{train\_gnn\_refine\_svd.py} を追加。素朴なGCNは埋め込みを崩して性能低下したため、学習開始時はベースライン同等（$\alpha=0$）になる残差混合に変更して安定化。
  \item 2026-03-18：データセット解析と訓練準備を拡張。\texttt{dataset\_report.txt} の変数TOP15に物理意味を併記し、ドメインを論文用に合併して図（ドメイン分布・埋め込み）を可読化。訓練ケースに \texttt{context\_paraphrased} を追加し、健全性チェック（ID整合・変数整合・重複）とベースライン評価を実装。さらに PyG の \texttt{equation\_graph.pt} を用いた retriever の1epoch動作確認スクリプトを追加（context+I/O 変数をクエリ特徴に含めて改善）。
  \item 運用ルール変更：全ての変動を開発ログに書き、変更のたびに \texttt{git commit} と \texttt{git push} で GitHub へ自動同期するよう Cursor ルールを強化（\texttt{.cursor/rules/development-log.mdc} および本ドキュメント技術メモを更新）。
  \item 開発ログに変動履歴を追記：モデル（gemini-2.5-pro 等）、デフォルト API キー、処理済み PDF 管理、\texttt{fetch\_papers.py}、\texttt{build\_graph\_data.py}、\texttt{generate\_training\_cases.py}（correct\_model\_ids の条件含む）、\texttt{analyze\_dataset.py}。GitHub 同期はコミット・push で行う。
  \item \texttt{build\_equation\_graph.py} を追加（PDF 一括、Handbook 生成、\texttt{unified\_equations.json} 統合）。処理済み PDF を \texttt{processed\_pdfs.json} で管理し未抽出のみ処理するよう変更。
  \item 開発記録の自動追記ルール（Cursor）とコミット時自動 push（post-commit フック）を導入。
  \item 初版：モデル変更、レートリミット対策、$c_w$/正規化の記録を追加。
  \item 2026-07-12（有意差検定スクリプト）：baseline vs reranker-10S の対応ある Wilcoxon 符号順位検定を実装（\texttt{analyze\_significance.py}）。設定A/B とも Recall@K\_correct・MAP・Recall@20 の全指標で p=0.00195（両側、10乱数全て reranker $>$ baseline のため理論最小値）。Cohen dz は Recall@K\_correct で11.9～12.9、MAP で9.4～32.3 と強効果。TDD に従い pytest で検証済み（tests/test\_significance.py）。出力は \texttt{experiments/significance\_stats.json}。
  \item 2026-07-12（greedy 3者比較の設定A拡張）：DAEのみ（設定B相当）で実施済みの3者比較（静的参照／推論のみgreedy／学習版greedy）を設定A（原型+複数文献+DAE、1,823件）に拡張。\texttt{run\_greedy\_3way\_A.sh} を新規作成し、reranker-10S・層化分割・top-k 200・10シード（計30ジョブ）で実行。集計は \texttt{analyze\_greedy\_3way.py} を引数化（\texttt{-{}-dir}／\texttt{-{}-out-json}／\texttt{-{}-out-fig}／\texttt{-{}-tag}。無引数時の既定挙動・出力はバイト同一で不変）して流用。結果（Recall@K\_correct、seed平均±SD、n=10）：静的参照 $0.7568 \pm 0.0560$、推論のみgreedy $0.7865 \pm 0.0487$、学習版greedy $0.7968 \pm 0.0530$ で、DAEのみと同じ 静的 $<$ 推論 $<$ 学習版 の順序を再現。検定（Wilcoxon 符号順位、両側）：train vs static は $\Delta=+0.0400$、p=0.00195 で有意、infer vs static は $\Delta=+0.0297$、p=0.00195 で有意。一方 train vs infer は $\Delta=+0.0103$、p=0.13086 で有意差なし（DAEのみでは p=0.03711 で有意だった差が、設定Aでは検出されず）。最難端 X$\ge$8 でも train vs static は $\Delta=+0.0481$、p=0.00195 で有意、train vs infer は $\Delta=+0.0086$、p=0.06445 で有意差なし。統計は \texttt{experiments/greedy\_3way\_A\_stats.json}、図は \texttt{docs/figures/fig\_greedy\_3way\_A.png}（生の per-case 出力 \texttt{experiments/xg\_A/} はコミット対象外）。
  \item 2026-07-12（修論Phase 1完了：全体検証とアウトラインPDFの確定）：\texttt{docs/superpowers/specs/2026-07-12-master-thesis-design.md} に基づく Phase 1（TeX骨格・章別3階層アウトライン・追加実験・英語図7点）が完了し、全体検証を実施。\texttt{thesis/master\_thesis/} で \texttt{latexmk -gg main.tex} が exit 0 で26ページの \texttt{main.pdf} を生成し、7章＋Appendix A--C を目次に確認、未解決参照（\texttt{??}）・未定義引用（\texttt{[?]}）はともに0件（\texttt{main.log} にも undefined 警告なし）。スペック照合：(1) 表紙（jsbook・和文表紙・英語本文）は研究室過去4本の形式を踏襲、(2) 全7章の P系列（トピックセンテンス）のみを通読し一貫した要約になることを確認、(3) 本文の主要数値を \texttt{experiments/strat\_A.json}・\texttt{strat\_B.json}・\texttt{significance\_stats.json}・\texttt{greedy\_3way\_stats.json}・\texttt{greedy\_3way\_A\_stats.json}・\texttt{feature\_x\_difficulty\_stats.json}・\texttt{dof\_stop\_stats.json} と1件ずつ突合し全て一致（記憶値なし）、Appendix A のハイパーパラメータも \texttt{strat\_A.json} の \texttt{config} と \texttt{set\_aware\_reranker.py} のデフォルト値に一致することを確認、(4) 第6章§6.4「$X\ge8$ で $+0.056$」は \texttt{greedy\_3way\_stats.json} の \texttt{tests\_hard\_X8plus.train\_vs\_static}（learned greedy vs 静的参照）を正しく参照しており、train vs infer との混同はないことを確認、(5) 追加実験1（有意差検定）・2（greedy3者比較の設定A拡張）は commit \texttt{8a8f4bf}・\texttt{e2914cc} で既にコミット済み。軽微な指摘：第4章§4.3（DAEケース生成の $X=1..10$／各100件）と§4.4（層化分割の6:2:2・候補400通り）の一部数値に直近の \texttt{\% source:} コメントが付いていない箇所があり、値自体は他章・他ファイルで裏付けられているため捏造ではないが、Phase 2 の本文化時に出典コメントを補うこと。Phase 1 完了につき、次は先生方への骨子（アウトライン）PDFレビュー依頼（ゼミ提出）。
  \item 2026-07-12（機構の考察一式と進捗報告）：加藤先生の「特徴量を設計して精度が上がっただけでは Kaggle と変わらない。なぜ上がるのかの考察が重要」とのご指摘に対応する 3 実験を実施し、進捗報告 \texttt{docs/progress\_report\_2026-07-12\_mechanism.tex/.pdf}（4 ページ）にまとめた。報告書では読みやすさを優先し、3 者比較（静的参照／推論のみ greedy／学習版 greedy）は「従来（基準固定）／逐次選択（学習・推論とも基準更新、＝学習版 greedy）」の 2 者に簡約して提示し（中間の「推論のみ greedy」と教師強制・scheduled sampling 等の専門用語は報告書からは省略。完全な 3 者比較は本記録と \texttt{greedy\_3way\_stats.json} に保持）、図は機構の層別図 1 点のみ掲載した。(1)\,特徴量×難易度クロス：\texttt{run\_feature\_x\_difficulty.sh}（reranker-7 に集合特徴量を 1 つずつ追加した 5 モード × 10 シード、設定 A・層化分割・per-case 保存）→ \texttt{analyze\_feature\_x\_difficulty.py}（シード対応 $t$ 検定＋正解式数・入力変数数の層別図 \texttt{fig\_mechanism\_fingerprint.png}、統計 \texttt{feature\_x\_difficulty\_stats.json}）。gComp $+0.052$（$p<0.001$）・gCoh $+0.050$（$p<0.001$）が上乗せのすべてを担い、gDom は $-0.000$（$p=0.98$）で寄与なし。上乗せは複数式・多変数のケースに限られ、入力 16 変数以上では $+0.035$ に縮小。前回報告の改善 $+0.222$ は「基本 7 特徴量の重み学習 $+0.172$」＋「集合特徴量 $+0.050$」に分解できることを確認。(2)\,学習版 greedy：\texttt{set\_aware\_reranker.py} に \texttt{-{}-train-greedy}（教師強制の逐次集合構築。各ステップで既選択集合を参照に set-aware 特徴量を再計算し、次の正解式を負例より上へ押す）と \texttt{-{}-seed-list}（並列実行用）を追加し、静的参照／推論のみ greedy／学習版 greedy の 3 者比較（DAE のみ・top-k 200・10 シード、\texttt{run\_greedy\_3way.sh}／\texttt{analyze\_greedy\_3way.py}）。Recall@K\_correct は $0.724 < 0.756 < 0.765$、train vs static $\Delta=+0.041$（$p<0.001$）、train vs infer $\Delta=+0.009$（$p=0.036$）、X$\ge$8 で差最大（$+0.056$）。実験 1 で予測した「固定参照による頭打ち」が介入で回復し、機構を因果として確認。(3)\,自由度ゼロ停止：\texttt{-{}-stop-dof}（\texttt{greedy\_close}／\texttt{set\_metrics}／\texttt{is\_closed} を新設）で、正解式数 $K$ を与えず「出力変数が全て被覆され、かつ未知数の数 $\le$ 式数」の時点で選択を停止（\texttt{run\_dof\_stop.sh}／\texttt{analyze\_dof\_stop.py}、統計 \texttt{dof\_stop\_stats.json}、図 \texttt{fig\_dof\_stop.png}）。完全一致率は $K$ 既知と同一（DAE 0.368／設定 A 0.526）、集合 F1 の低下は僅少（設定 A $-0.008$、$p=0.006$）、解ける系の割合は 0.88／0.94 で上位 $K$ 件の 0.78／0.87 を上回る。事前検証：正解集合の自由度ゼロは全 variant で成立（オリジナル 92・文献横断 731・DAE 1{,}000 の全件）、部分集合での早期成立なし（DAE 1{,}000 件 × 8 順序で 0 件）。ただし $K$＝出力変数数がほぼ全件で成立するため $K$ 予測自体は自明である点を報告に明記。コード・図・統計は commit \texttt{08c371d}。進捗報告の図 3 点は作文ルールに合わせ文言を清書（比喩・プログラム変数名・未定義略語を排除）。
  \item 2026-07-12（baseline$\to$7 特徴量の分解）：「集合特徴量の内訳は説明したが、baseline（古典的検索）から 7 特徴量 MLP への $+0.172$ の内訳が未説明」との本人指摘を受け、中間条件 \texttt{reranker-2}（baseline と同じ 2 特徴量〔text\_sim・io\_jaccard〕のみを MLP で学習。\texttt{MODE\_CONFIGS} に \texttt{base\_subset} を追加、既存モード非破壊）を実装し、設定 A・層化分割・10 シードで実行（\texttt{run\_reranker2.sh}、出力 \texttt{experiments/xd/reranker-2\_\_*.json}）。結果（Recall@K\_correct、シード対応 $t$ 検定）：baseline $0.521 \to$ 同 2 特徴量 MLP $0.601$（\textbf{重み付けの学習 $+0.080$}、$p<10^{-6}$）$\to$ 7 特徴量 $0.693$（\textbf{特徴量の追加 $+0.092$}、$p<10^{-6}$）$\to$ 集合特徴量 $0.743$（$+0.050$）。合計 $+0.222$ で前回報告と一致し、\textbf{改善の全 3 段が分解できた}。進捗報告（4 ページ版）の \S2・表 1・\S3・考察・まとめに反映。あわせて本人の手動修正の誤字等を修正（タイトル「訓練特徴量」$\to$「設計した特徴量」、ご指摘引用の「効くの考察」$\to$「効くのかの考察」、\S1 の「前回」重複）。追加した 5 特徴量（意味類似・入力被覆・出力被覆・式の特化度・分野一致）のどれが $+0.092$ を担うかは未分解（報告書の限界に明記）。
  \item 2026-07-14（進捗報告を実験 1 のみに絞る）：本人の指示により、進捗報告 \texttt{progress\_report\_2026-07-12\_mechanism} から実験 2（逐次選択）・実験 3（自由度ゼロ停止）の節を本文ごと削除し、\textbf{実験 1（精度向上の分解と層別）のみの 3 ページ構成}に変更。タイトル後半を「精度向上の分解と、効く場所の特定」に、考察を 2 項目（課題の性質・限界）に、まとめを分解の結果のみに整合。その後、本人の指示（「次回にも書く必要ない、進捗報告でしない」）により今後の課題からも逐次選択・停止規則の予告を削除し、報告書から実験 2・3 への言及を完全になくした（限界 (iii) も「推定である」の表現に整理）。実験 2・3 の完全な結果は本記録（2026-07-12 の項）と \texttt{greedy\_3way\_stats.json}／\texttt{dof\_stop\_stats.json} に保持。
  \item 2026-07-14（全特徴量の個別分解：$+0.092$ の内訳と話題特徴の無効性）：本人指摘「『上乗せ』の比較先が曖昧」「『どの特徴量が効くか』の実験で基本 7 特徴量を対象から除いている」を受け、基本特徴量も 1 つずつ分解した。\texttt{MODE\_CONFIGS} に \texttt{reranker-2+Svd/InCov/OutCov/Spec/Dom}（土台＝2 特徴量に基本特徴量を 1 個ずつ追加。\texttt{base\_subset} で列を選択）を追加し、設定 A・層化分割・10 シードで実行（\texttt{run\_base\_ablation.sh}）。結果（土台 A＝2 特徴量 MLP $0.601$ からの差、シード対応 $t$ 検定）：\textbf{式の特化度 $+0.090$（$p<0.001$）が $+0.092$ のほぼ全部を担い、これ 1 個で 7 特徴量（$0.693$）にほぼ到達}。入力被覆 $+0.019$（$p<0.001$）・出力被覆 $+0.006$（$p=0.002$）は小、\textbf{意味の近さ（SVD）$-0.002$（$p=0.48$）・分野の一致 $+0.002$（$p=0.06$）は効果なし}。集合版の分野一致（gDom $p=0.98$）と合わせ、\textbf{話題の近さを見る特徴量はどの形でも効かない}という対称的な結論を得た。進捗報告の表 1 を全 10 特徴量の個別効果（土台 A/B の 2 段はしご）に再構成し、「上乗せ」を「土台からの差」に統一（比較先を明示）、結論・考察・まとめを「効くのは変数の重なり、効かないのは話題の近さ」に書き換え。3 ページ維持。追加 5 特徴量の内訳は分解済みとなり、限界からは削除。統計は per-case JSON（\texttt{experiments/xd/reranker-2\_*\_\_*.json}）に保持。
  \item 2026-07-14（分野一致度の負の結果を修論本文に明文化）：本人の問い「分野一致度の特徴量が精度向上に効かないのは、分野分けが広すぎて個別化できないためか。分野をさらに細かくすべきか」を実データで検証し、前提が逆であることを確認した。特徴量が照合する \texttt{domain} フィールドは、論文収集に用いた 32 分野（\texttt{fetch\_papers\_32domains.py}）ではなく、\textbf{抽出時に大規模言語モデル（Gemini 2.5 Pro）が式ごとに自由記述で付けた分野名}である。\texttt{unified\_equations.json}（11{,}146 式）で集計すると、異なる分野ラベルは \textbf{1{,}998 種}、うち 27.9\% が 1 回のみの出現、1 論文あたり平均 6.36 種で、表記ゆれ（``Transport Phenomena'' と ``Transport phenomena'' を別物として扱う）も多い。したがって分野が広すぎるのではなく細かすぎ、完全一致（集合特徴量）も部分文字列一致（基本特徴量）もまばらにしか成立しないため信号にならない。分野をさらに細かくすると疎性が悪化するだけであり、話題の近さは文章類似度・変数の一致度がすでに担う（本課題は同じ話題ではなく同じ変数で解ける）。数値照合：基本版の分野一致は土台 2 特徴量 $0.6010 \to 0.6034$（$+0.0024$、対応 $t$ 検定 $p=0.063$、Wilcoxon $p=0.123$、\texttt{experiments/xd/reranker-2\_\_*.json}・\texttt{reranker-2\_Dom\_\_*.json} の 10 シード）、集合版の分野一致（gDom）は $-0.0001$（$p=0.98$、\texttt{feature\_x\_difficulty\_stats.json}）で、いずれも効果なし。この負の結果を修論本文 \texttt{thesis/master\_thesis/ResultsAndDiscussion.tex} に独立節 \S6.6「The Null Result for Topic Agreement」として追加し（Limitations の直前。既存 Limitations は \S6.7 に繰り下げ、P7$\to$P8）、分野ラベルの生成主体を第 4.1 節（Gemini 2.5 Pro）と整合させた。\texttt{uplatex} で目次・本文の描画を確認済み。
  \item 2026-07-15（分野ラベル合併の検証：合併しても domain 特徴は効かない）：\S6.6 の「分野一致度が効かないのは分野が細かすぎるため」に対し、「ラベルを合併して粗くすれば効くようになるか」を検証した。コミット済み \texttt{set\_aware\_reranker.py} の \texttt{run\_mode} を外部から呼び、式ごとの分野ラベル列（\texttt{eq\_domains}）だけを合併ラベルに差し替えて基本版 domain 特徴つき \texttt{reranker-2+Dom}（基本2特徴量＋分野一致）を再学習する検証スクリプト \texttt{run\_domain\_merge\_ablation.py} を新規作成（パイプライン本体は無改変）。合併は3粒度：L1（小文字化・記号除去・語尾複数形の統合で 1{,}998$\to$1{,}958 群）、KMeans による 100 群、KMeans による 40 群（32 収集分野スケール）。短い分野名では TF-IDF＋average-linkage 凝集クラスタリングが単一クラスタに退化したため、L2 正規化 TF-IDF 上の KMeans を採用。設定 A・層化分割・top-k 50・10 seed で、domain 不使用の \texttt{reranker-2}（$0.6010$、既存 \texttt{experiments/xd/} を再利用）を base とし対応 $t$ 検定：\textbf{Recall@K\_correct の差はどの粒度でも $\approx 0$ で全て非有意}（orig $+0.0024$／$p=0.06$、L1 $+0.0006$／$p=0.68$、K100 $-0.0001$／$p=0.84$、K40 $+0.0004$／$p=0.63$）。元ラベルのわずかな $+0.002$ すら合併で消える。\textbf{障害は疎性でなく冗長性}：第1段の候補集合が既に話題的に近いため、合併した分野一致は正解候補と不正解候補を分離しない。この結果で修論 \S6.6 を「分野を細かくすると疎性が悪化する（推定）」から「細かくしても疎性が悪化し、粗くしても（合併しても）寄与は $\approx 0$（実測）」の両方向の主張に格上げした。統計は \texttt{experiments/domain\_merge\_stats.json}。

  \item 2026-07-21（P-77 紹介スライド）：PSE ASIA 2026 のポスター P-77（台湾科技大、LLM＋MCP による Aspen Plus 操作自動化）を研究室ゼミで紹介するスライド 8 枚を作成。仕様 \texttt{docs/superpowers/specs/2026-07-21-p77-seminar-slides-design.md}・計画 \texttt{docs/superpowers/plans/2026-07-21-p77-seminar-slides.md} に基づき、本文・台本を \texttt{scripts/seminar\_p77/content\_p77.py}、描画を \texttt{scripts/seminar\_p77/build\_p77\_slides.py} に実装（PSE Asia 発表テンプレのマスターを流用、自作図 3 点、ノート欄に台本約 12 分）。pytest（\texttt{tests/test\_p77\_content.py}・\texttt{tests/test\_p77\_slides.py}）で 8 枚構成・全テキスト 20pt 以上・略称の正式名称併記・禁止語（格上げ表現含む）・台本分量を機械検証。出力 \texttt{\textasciitilde/Desktop/ゼミ/論文紹介\_P77\_LLM\_Aspen\_かずひろ.pptx} はリポジトリ外のためコミット対象外。
  \item 2026-07-21（紹介スライド v2 改稿）：発表者レビュー（「何をやっているか分かりにくい」「抄録だけでは体系的な話に足りない」）を受け、ゼミ紹介デッキを P-77 単独 8 枚から\textbf{DTU 論文中心の 10 枚}に全面改訂（発表 12 分は維持）。中心例を全文公開の Liang, Groll \& Sin, arXiv:2601.11650（Claude Desktop + 自作 MCP サーバ + AVEVA Process Simulation、CC BY 4.0）に変更し、「向く作業・やるべき作業・任せてはいけないこと」の 3 つの問いを背骨に再構成。ツール数は HTML の Table 1 を直接数えて 16 個と確認（抽出モデルは 18/19 とぶれたため一次確認を実施）。掲載数値はすべて論文実測（ツール 7 回で約 2{,}000 変数から要点 24 個、改善案 11 個の著者採点 3/4/3/1/0、還流比反復で純度 84.2$\to$95.1 mol\%、算数ミス $+10.9\%\leftrightarrow+12.9\%$、省エネ主張 20〜40\% 対 専門家見積 2〜5\%）。実装は \texttt{content\_p77.py} 全面書き換え・\texttt{build\_p77\_slides.py} 改修（表の可変幅/高さ、構成図 5 箱、回避ルート図削除、出力名・docProps 更新）、pytest を 22 本に拡充（10 枚構成・論文数値一致・研究室 2 件言及の検証を追加）。出力 \texttt{\textasciitilde/Desktop/ゼミ/論文紹介\_LLM\_MCP\_シミュレータ操作\_かずひろ.pptx}（v1 の P-77 版ファイルは残置）。
  \item 2026-07-21（スライドのレイアウト刷新）：本人指摘（スライド 1 で本文が注記に重なる等）を受け、固定座標レイアウトを廃止し\textbf{実フォント計測にもとづく動的配置}に全面刷新。\texttt{scripts/seminar\_p77/textmetrics.py} を新設し、実際に使うフォント（欧文 IBM Plex Sans／和文ヒラギノ角ゴシック）のグリフ幅を matplotlib 経由で計測して折り返し行数を算出、全レンダラを「前要素の実測下端から次要素の y を決める」積み上げ方式に変更（表は行高もセル内折り返しから実測）。あわせてレイアウト監査 \texttt{audit\_presentation}（下端はみ出し・図形内あふれ・テキスト矩形の重なりを検査）と \texttt{tests/test\_p77\_layout.py} を追加。監査の検出力は旧ビルドに適用して重なり 5 件（S2 リード×表、S2/S5/S7 注記×出典、S4 戻りラベル×箇条書き）を検出することで確認し、新ビルドは違反 0 件。全 23 テスト green。
  \item 2026-08-09（話題特徴が寄与しない原因の実証と正解・不正解ケース分析）：加藤先生の宿題「性能向上手法として期待したことと結果の関連と原因の考察（例：精度向上にほぼ貢献できなかった特徴量）。正解・不正解のケースを見て分析。意味の近さが役立たない原因は TF-IDF$\to$SVD のベクトルがケース間で差が小さかった可能性」に対応する 2 分析を実施し、\texttt{docs/考察\_特徴量の期待と結果\_2026-08-09.md} にまとめた。(1)\,診断（モデル非依存、\texttt{diagnose\_topic\_features.py}・pytest 11 本）：\textbf{仮説は支持されない}——ケース間クエリベクトルのコサイン類似は中央値 0.218（SVD 空間）・0.070（生 TF-IDF）で 0.5 超は 5.4\% にとどまり、svd\_sim 単独の全 DB AUC は $0.950\pm0.003$（ベクトルが縮退していれば不可能）。真因は候補絞り込み後の冗長性で、stage1 候補 50 件内では svd\_sim の AUC が 0.594 に崩壊（specificity 0.965、変数系 0.72--0.97 は保持）。さらに stage1（0.7$\times$text\_sim+0.3$\times$jaccard）の混合選抜が候補内に $\mathrm{corr}(\text{text\_sim},\text{io\_jaccard})=-0.404$ の見かけの負相関を作り、text\_sim は候補内 AUC 0.437 と逆向きになることを実測。被覆率（正解式が候補 50 件に入る割合）は混合 $0.920\pm0.004<$ Jaccard 単独 $0.961\pm0.002$ で、stage1 の混合重みが上限を下げていることも判明（次の介入実験候補）。(2)\,正解・不正解ケース分析（\texttt{analyze\_case\_outcomes.py}・seed=42・テスト 423 件・pytest 5 本）：取り逃した正解式 435 本の \textbf{61.8\% は stage1 見逃し}、38.2\% が再ランク見逃し。落ちたケースは正解式数 5.3 vs 2.1・入力変数 17.8 vs 10.8（いずれも Mann-Whitney $p<0.001$）で既報の難易度プロファイルと一致。反実仮想の救済可能性 P(見逃し正解 $>$ 誤上位) は svd\_sim $0.516\pm0.040$・domain 0.494 とコイン投げ同水準（io\_jaccard は 0.375 で誤り側が優位）＝\textbf{話題特徴を強めても残存誤りは救えない}。具体例で誤りの型を同定：別文献の数学的同一式（話題類似はむしろ誤りを推す。$dX/dt=\mu X$ の重複）、式の向き（どの変数を左辺で決めるか）の欠落、類似式ファミリによる候補占拠。出所検証では誤り候補の同一文献率は 22.2\% に過ぎず、話題の均質さは文献重複＋選抜由来。図 \texttt{docs/figures/fig\_topic\_diagnosis.png/.pdf}（(a) ケース間類似分布、(b) 候補内 AUC・SEM エラーバー付き）、統計 \texttt{experiments/topic\_feature\_diagnosis.json}・\texttt{case\_outcomes\_seed42.json}。並行セッションが同ブランチで PFI per-case 本実行（\texttt{pfi\_per\_case.json} 生成、約 2 時間）を担当中のため、本セッションからの重複実行は停止して委譲。
  \item 2026-08-09（修論 第3--5章の本文化・ケース族の用語改訂・第6章に誤り分析節を新設）：本人指示の今後の方針「現状のデータセット・提案手法・評価方法を論文形式（Methods / Experiments 相当）でまとめる。オリジナル・文献横断・DAE の言葉の使い方は変更する。論文の主張を考える」に対応。(1)\,\textbf{用語改訂}：正解集合の作り方という単一の軸で命名し直した（\texttt{original}$\to$\textbf{as-published}／文献掲載、\texttt{multisource\_*}$\to$\textbf{cross-document}／文献横断、\texttt{dae\_X*}$\to$\textbf{synthesized}／合成、拡張$\to$augmented）。旧称は出所と作り方の 2 軸を混ぜており、「オリジナル」はケースの性質でなくデータセットの履歴を指し、「DAE」は他の族も満たす性質（文献掲載・文献横断のケースも微分代数系）を族名にしていた。設定 A/B は「混合」「合成のみ」と併記。全章（Introduction / Dataset / Experiment / Results / Conclusion / Appendix）に適用。(2)\,\textbf{第3章 Method の本文化}：問題定義（自由度 $\mathrm{DoF}(S)=|\text{未知数}|-|S|$ と閉包条件）・第1段・第2段・逐次選択・自己停止を段落形式で執筆し、特徴量表（10 特徴の定義と群）を新設。\textbf{実装を一次情報として骨子の誤記 4 件を訂正}：①第1段のスコアは 2 特徴の混合（$0.7\,\text{text\_sim}+0.3\,\text{io\_jaccard}$）であり「7 特徴」ではない（7 特徴は第2段 MLP の入力）、②MLP の隠れ層は 64 ユニット\textbf{2 層}（骨子は 1 層）、③集合特徴の参照集合は第1段上位 \texttt{K\_REF}=5 件（候補集合 $k=50$ とは別物）、④閉包条件は $\mathrm{DoF}\le 0$ かつ出力被覆（骨子の $\mathrm{DoF}=0$ は等号のみで狭い）。(3)\,\textbf{第4章 Dataset の本文化}：一次データから全数値を再実測（式 11,146＝PDF 328 文献 10,437 式＋生成参照集合 33 種 709 式、変数記号 13,220、分野ラベル 1,998・singleton 27.9\%・文献あたり 6.36）。族ごとの出所を実測（文献横断は 100\% が 2 文献以上・最大 3、文献掲載 1.1\%・合成 4.1\%）、評価 1,823 件の 90.8\% が 2 式以上・平均 3.90 式。\textbf{K＝出力変数数が文献横断・合成で 100\%（文献掲載 92.4\%）}＝K 予測が自明であることを本文に明記。(4)\,\textbf{第5章 Experiment の本文化}：主指標 Recall@K\_correct が\textbf{R-precision と同一}であること（窓＝正解数なので precision と recall が一致、\texttt{evaluate\_multi\_eq.py} で確認）、拡張除外の根拠（含めると baseline 0.521$\to$0.622 で差が $+0.222\to+0.156$ に縮む）、\textbf{候補窓のハンディキャップ}（baseline は全 DB を順位付けするが再ランクは候補被覆 0.920 が上限）、\textbf{LLM 比較の生成器と判定器が同一モデル}（ともに claude-opus-4-8）で LLM 側に有利な保守的比較であることを明記。(5)\,\textbf{候補窓の被覆率を実測}：$k=50$ で 0.9203（全正解被覆ケース 77.2\%）、$k=100$ 0.9569、$k=200$ 0.9775（90.3\%）、$k=400$ 0.9895。(6)\,\textbf{第6章}：§6.6 を「分野一致」から「話題全般（Topic Similarity）」に拡張し P7$'$（ベクトル差仮説の棄却・全 DB AUC 0.950 対 候補内 0.594・選択効果 $-0.404$）を追加、\textbf{§6.7 Error Analysis を新設}（誤りの 61.8\% は第1段・被覆 0.920 対 変数単独 0.961・救済可能性 0.516・等価式重複と式の向きの 2 型）。既存 Limitations は §6.8 に繰り下げ、節参照を \texttt{\textbackslash label}/\texttt{\textbackslash ref} 化。(7)\,\textbf{主張の文書化}：\texttt{docs/論文の主張\_2026-08-09.md} に貢献 C1--C7・主張しないこと 5 項（自動構築した／K を当てられる／LLM より優れた手法である 等）・想定質問 Q1--Q5・用語対応表をまとめた。ビルド確認：\texttt{latexmk -gg main.tex} exit 0、21 ページ、未解決参照・未定義引用なし。
  \item 2026-08-09（進捗報告：話題の近さの限界と誤りの所在）：今回の議論と追加実験を進捗報告 \texttt{docs/progress\_report\_2026-08-09\_topic.tex/.pdf}（5 ページ）にまとめた。構成は決定版見本（6\_17 進捗報告）に準拠し、\S1 前回ゼミレビュー、\S2 今回の進捗（番号リスト 6 項目）、\S3 実験 1（意味の近さはどこで力を失うか）、\S4 実験 2（正解・不正解のケースを見る）、\S5 修士論文の本文化とケースの呼び方の整理、\S6 考察、\S7 まとめと今後の課題。加藤先生のご指摘（期待と結果の関連・原因の考察／正解不正解ケースの分析／TF-IDF$\to$SVD のベクトル差仮説）を \S2 冒頭に明示し、以降がその回答になる形にした。\textbf{新規の追加実験}：第 1 段の絞り込み方ごとの被覆率を $k=10,25,50,100,200,400$ で測る \texttt{analyze\_stage1\_coverage.py}（pytest 6 本、\texttt{-{}-figure-only} で再計測なしの作図が可能）を新規作成。現行の混合（文章 0.7＋変数 0.3）は $k=50$ で 0.920・$k=200$ で 0.977 に対し、\textbf{変数の一致度だけで絞ると $k=50$ で 0.961}、文章類似度だけでは 0.431、意味の近さだけでは 0.518 であった（\texttt{experiments/stage1\_coverage.json}）。被覆率は第 2 段 Recall@K\_correct の上限を与えるので、第 1 段を変数寄りに再設計する介入実験の根拠になる。図は新規 \texttt{docs/figures/fig\_stage1\_coverage.png}（(a) 絞り込み方別の被覆率、(b) 救済可能性。ともに誤差棒＝標準誤差）。\textbf{図の清書}：\texttt{fig\_topic\_diagnosis} のパネル (b) が \texttt{specificity}・\texttt{io\_jaccard} 等のプログラム変数名のままで作文ルール A6（アンダースコア表記は文書・論文で使わない）に反していたため、\texttt{diagnose\_topic\_features.py} に日本語ラベル表（式の特化度・変数の一致度・意味の近さ 等）を追加して再生成し、軸ラベルの ``SEM'' も「標準誤差」に、``stage1'' も「第 1 段」に改めた。再実行で \texttt{topic\_feature\_diagnosis.json} はバイト単位で不変（決定的であることを確認）。
  \item 2026-08-09（コサイン類似度の物差し）：本人の「コサイン類似度の中央値 $0.218$ のイメージがつかない。$0.218$ ってどのくらい違うのか」という問いに答えるため、同じ空間の中に基準となる組を置いて較正した（\texttt{calibrate\_similarity\_scale.py}、pytest 8 本）。式本文で学習した TF-IDF$\to$SVD 空間へ全 2{,}838 ケースを射影し、コサイン類似度の中央値を組の種類別に測ると、\textbf{同じ核モデルから入出力だけ変えた組 $1.000$}（四分位 $0.982$–$1.000$、n=1{,}050）、\textbf{同じケースの説明文を書き換えただけの言い換えの組 $0.944$}（$0.926$–$0.962$、n=276）、\textbf{同じ文献に由来する別ケースの組 $0.445$}（$0.302$–$0.588$、n=16{,}123）、\textbf{無作為な 2 ケースの組 $0.218$}（$0.149$–$0.313$、n=199{,}882）であった。決め手は\textbf{言い換えの組 276 組のうち $0.218$ 以下が 1 組もない}（0.0 パーセンタイル）ことで、内容が同じ組と無作為な組は分布が重ならない。$0.218$ 付近の実例は「微生物がグルコースを消費する発酵系」と「PI 制御下のタンク液位の応答」で話題は明確に異なる。全ケースが英語の化学工学の文章のため共通語彙による下駄があり、無関係でも $0$ にはならない（実質的な範囲は $0.15$〜$1.0$）。この較正により \S3 の「ベクトル差が小さいという仮説の棄却」が目に見える形になったので、進捗報告 \S3 に 1 項目と箱ひげ図（\texttt{docs/figures/fig\_similarity\_scale.png}、図 2）を追加し 6 ページとした。統計は \texttt{experiments/similarity\_scale.json}。
  \item 2026-08-09（進捗報告の改訂：PFI の追加と 4 ページへの圧縮）：本人の指摘を受けて進捗報告 \texttt{docs/progress\_report\_2026-08-09\_topic.tex} を段階的に改訂した。(1)\,\textbf{候補件数の取り違えを訂正}：「データ拡張後に第 1 段の候補を 200 件に上げた」との指摘で確認したところ、主要比較 \texttt{strat\_A/B} は $k=50$、逐次選択・自己停止（2026-07）は $k=200$ と使い分けていた。今回の分析は全て $k=50$ で行っていたため $k=200$ で測り直し（\texttt{analyze\_case\_outcomes.py}・\texttt{diagnose\_topic\_features.py} に \texttt{-{}-top-k} を追加、既定 50 で挙動不変）、\textbf{「誤りの主因は第 1 段（61.8\%）」が $k=50$ 限定の主張}であったことが判明した。$k=200$ では第 1 段由来は 28.4\% に下がり並べ替えの誤りが 71.6\% と逆転する（見逃し 435 本$\to$373 本、Recall@$K$ 0.791$\to$0.807）。一方で「話題の近さでは救えない」は強まった（救済可能性 0.516$\to$0.421 と 0.5 を下回る）。(2)\,\textbf{PFI が未報告であったことの訂正}：前回報告書（ヘッダ 2026-07-14）に PFI の図が入っているのを見て「前回報告済み」と書いたが、追記コミットは \texttt{cefe85a}（2026-07-15）＝ゼミ後で、配布版には入っていなかった。PFI は一度も報告していないため、動機・方法・全体結果（変数群 $0.743\to0.065$、話題群 $+0.006$ n.s.、特化度 $+0.617$）から書き起こす節に改めた。あわせてケース別分解の結果——\textbf{分野の一致の寄与ゼロは「効かない」ではなく相殺}（正しい向きのケース 1{,}169 件で $+0.012$、逆向き 414 件で $-0.034$＝シャッフルすると精度が上がる、$p=1.3\times10^{-8}$）——を追加した。(3)\,\textbf{読み手に合わせた書き直し}：チーム内語（「変数寄り」「文章 0.7＋変数 0.3」）を平易化し、§1 に 2 段構えの説明を追加（それが無いまま「第 1 段」「候補 50 件」を使っていた）。前方参照の補足（「後述の AUC で 0.950」）と、主語のない「取り逃した正解式」を排除。修論に関する記述（用語整理の節）は進捗報告から全削除。定義のない造語「合成ケース」を「設定 B（DAE のみ 1{,}000 件）」に戻した。(4)\,\textbf{4 ページに圧縮}：本人指示により 7 ページ$\to$4 ページ。考察節をまるごと削除し、難しい議論（選択効果＝候補内の見かけの逆相関、コサイン類似度の較正、依存ケースの属性分析）と図 2 点を落とした。残したのは仮説の検証・候補を絞った後で効かなくなること・誤りの所在・誤りの 2 型・PFI と相殺。比喩（「下駄」「足を引っ張る」「居座る」「コイン投げ」）を全て平易な表現に置換。新規スクリプト \texttt{analyze\_stage1\_coverage.py}（pytest 6 本）・\texttt{calibrate\_similarity\_scale.py}（pytest 8 本）は残し、統計は \texttt{experiments/} に保持（報告書には非掲載）。
\end{itemize}

\end{document}
